% macro conflict, first definition kept: \newcommand{\Hom}{\operatorname{Hom}}
% macro conflict, first definition kept: \newcommand{\uHom}{\underline{\operatorname{Hom}}}
% macro conflict, first definition kept: \newcommand{\Ob}{\operatorname{Ob}}
% macro conflict, first definition kept: \newcommand{\Ob}{\operatorname{Ob}}
% macro conflict, first definition kept: \newcommand{\Hom}{\underline{\mathrm{Hom}}}
% macro conflict, first definition kept: \newcommand{\Fl}{\mathrm{Fl}}
% macro conflict, first definition kept: \newcommand{\Hom}{\underline{\mathrm{Hom}}}
% macro conflict, first definition kept: \newcommand{\Hom}{\operatorname{Hom}}
% macro conflict, first definition kept: \newcommand{\Ob}{\mathrm{Ob}}
% macro conflict, first definition kept: \newcommand{\Ch}{\mathrm{Ch}}
% macro conflict, first definition kept: \newcommand{\Ch}{\mathrm{Ch}}
\documentclass[11pt]{article}
\usepackage{fontspec}
\usepackage[french]{babel}
\usepackage{amsmath}
\usepackage{amssymb}
\usepackage{amsthm}
\usepackage{mathrsfs}
\usepackage{tikz-cd}
\usepackage{enumitem}

\newcommand{\printedpage}[1]{\marginpar{\footnotesize p.~#1}}
\newcommand{\msp}[1]{\marginpar{\footnotesize [#1]}}
\newcommand{\Hot}{\mathrm{Hot}}
\newcommand{\D}{\mathbf{D}}
\newcommand{\uHom}{\underline{\mathrm{Hom}}}
\newcommand{\hm}{\mathrm{hom}}
\newcommand{\Hom}{\mathrm{Hom}}
\newcommand{\homext}{\mathrm{hom}\,\mathrm{ext}}
\newcommand{\HOT}{\mathrm{HOT}}
\newcommand{\Ob}{\mathrm{Ob}\,}
\newcommand{\Ltens}{\overset{\mathbf{L}}{\otimes}}
\newcommand{\RHom}{\mathbf{R}\mathrm{Hom}}
\newcommand{\ab}{\mathrm{ab}}
\newcommand{\defeq}{\overset{\text{déf}}{=}}
\newcommand{\cA}{\mathcal{A}}
\newcommand{\bDelta}{\boldsymbol{\Delta}}
\newcommand{\id}{\mathrm{id}}
\newcommand{\Fl}{\operatorname{Fl}}
\newcommand{\uCh}{\underline{\mathrm{Ch}}}
\newcommand{\bPhi}{\boldsymbol{\Phi}}
\newcommand{\uFl}{\underline{\mathrm{Fl}}}
\newcommand{\Ch}{\underline{\mathrm{Ch}}}
\newcommand{\card}{\mathrm{card}}
\newcommand{\cM}{\mathcal{M}}
\newcommand{\Chh}{\underline{\underline{\mathrm{Ch}}}}
\newcommand{\KK}{\underline{\underline{K}}}
\newcommand{\dfeq}{\overset{\text{déf}}{=}}
\newcommand{\Cathot}{\mathrm{Cat\,hot}}
\newcommand{\Ind}{\operatorname{Ind}}
\newcommand{\Ens}{\mathrm{Ens}}
\newcommand{\Cat}{\mathrm{Cat}}
\newcommand{\uuK}{\underline{\underline{K}}}
\begin{document}
% ---- from Der-ChIII-567_p001-005.tex
% ===================== page 1 (title page) =====================
\printedpage{--}
% no page number printed on this page
% head: (none)

\begin{center}
{\LARGE\sffamily LES DÉRIVATEURS}

\bigskip

\textbf{ALEXANDRE GROTHENDIECK}

\bigskip\bigskip

\textbf{Chapitre III}

\medskip

\textbf{Hom externes dans les dérivateurs}
\end{center}

\clearpage

% ===================== page 2 (editorial note) =====================
\printedpage{--}
% no page number printed on this page
% head: (none)

Ce texte a été déchiffré et transcrit en \LaTeXe{} par M. Künzer. Il a été édité par M. Künzer, J. Malgoire, et G. Maltsiniotis. La transcription est aussi fidèle que possible au manuscrit. Pour les quelques corrections évidentes, ou rares commentaires des éditeurs, ainsi que pour la numérotation originale des pages du manuscrit, les caractères de machine à écrire \texttt{[typewriter]} entre crochets sont utilisés. Un point d'interrogation entre crochets signifie que l'on n'est pas sûr du mot qui précède.

Cette édition est provisoire. Les remarques, commentaires et corrections sont bienvenus. Envoyer un message à :

\textbf{[e-mail address removed]}
\begin{flushright}
G. Maltsiniotis\qquad\qquad
\end{flushright}

\clearpage

% ===================== page 3 (printed 1) =====================
\printedpage{1}
% head: A. Grothendieck : Dérivateurs, Ch. III, Édition des Universités de Montpellier II et Paris VII
% foot: Version 2 février 2002

\msp{page 1}

\section*{1\quad Relation des dérivateurs avec la catégorie $\Hot$}

Soit $\D$ un dérivateur, $\mathcal{A}$ sa catégorie fondamentale,
\[
\mathcal{A} = \D(e).
\]
Je me propose de définir, si faire se peut, des foncteurs
\[
\left\{
\begin{array}{rccccccl}
\Hot & \times & \mathcal{A} & \xrightarrow{\pi_{\D}} & \mathcal{A} & (\xi, M) & \longmapsto & \xi \times M \ \text{ ou } \ \xi \otimes M \quad \text{(suivant les cas)}\\
\Hot^o & \times & \mathcal{A} & \xrightarrow{\varpi_{\D}} & \mathcal{A} & (\xi, M) & \longmapsto & \hm(\xi, M).
\end{array}
\right.
\tag*{(1)}
\]
Ces foncteurs doivent s'échanger par dualité : si on prend les foncteurs $\pi^o$, $\varpi^o$ sur les catégories opposées, on trouve
\[
\left\{
\begin{array}{rcccc}
\Hot & \times & \mathcal{A}^o & \xrightarrow{\varpi^o_{\D}} & \mathcal{A}^o\\
\Hot^o & \times & \mathcal{A}^o & \xrightarrow{\pi^o_{\D}} & \mathcal{A}^o,
\end{array}
\right.
\tag*{(1$'$)}
\]
et on doit avoir
\[
\left\{
\begin{array}{rcl}
\varpi^o_{\D} & = & \pi_{\D^o}\\
\pi^o_{\D} & = & \varpi_{\D^o}.
\end{array}
\right.
\]
À vrai dire, considérant (1) comme des foncteurs
\[
\begin{array}{rcl}
\Hot & \xrightarrow{\pi_{\D}} & \uHom(\mathcal{A}, \mathcal{A})\\
\Hot^o & \xrightarrow{\varpi_{\D}} & \uHom(\mathcal{A}, \mathcal{A}),
\end{array}
\]
on s'attend à ce que ces foncteurs s'étendent en
\[
\left\{
\begin{array}{rcl}
\Hot & \xrightarrow{\pi_{\D}} & \uHom(\D, \D)\\
\Hot^o & \xrightarrow{\varpi_{\D}} & \uHom(\mathcal{A}, \mathcal{A})
\end{array}
\right.
\tag*{(2)}
\]
\texttt{[plutôt deux fois} $\uHom(\D, \D)$\texttt{]}. Cela signifie, essentiellement, que pour $\xi \in \mathrm{Ob}\,\Hot$ donné, les opérations $\xi \times_{\D} M$ et $\uHom_{\D}(\xi, M)$ \texttt{[plutôt} $\hm_{\D}(\xi, M)$\texttt{]} sont définies non seulement dans $\mathcal{A} = \D(e)$, mais dans tous les $\D(I)$, et ceci de façon à commuter aux foncteurs $f^*$ :

\msp{page 2}
Si
\[
f : I' \longrightarrow I,
\]
\printedpage{2}
% head: A. Grothendieck : Dérivateurs, Ch. III, Édition des Universités de Montpellier II et Paris VII
% foot: Version 2 février 2002
alors pour $F \in \mathrm{Ob}\,\D(I)$, on doit avoir des isomorphismes canoniques
\[
\left\{
\begin{array}{rcl}
f^*(\xi \times_{\D} M) & \simeq & \xi \times_{\D} f^*(M)\\
f^*\hm_{\D}(\xi, M) & \simeq & \hm_{\D}(\xi, f^*(M)),
\end{array}
\right.
\tag*{(3)}
\]
avec les compatibilités habituelles pour un composé $fg$, relatif à $I'' \xrightarrow{g} I' \xrightarrow{f} I$.

Mais je me place de préférence dans $\mathcal{A}$, ce qui devrait impliquer les constructions similaires dans tous les $\D(I)$, quitte à appliquer ce qui a été dit au dérivateur $\D_I$.

L'intuition algébrique fondamentale est celle-ci, en désignant par
\[
X \longmapsto \mathrm{hot}(X) : (\mathrm{Cat}) \longrightarrow (\Hot)
\]
le foncteur de localisation de (Cat) dans (Hot), obtenu en inversant les flèches dans (Cat) qui sont des `équivalences faibles'. Supposons $X$ dans Diag, soit
\[
p : X \longrightarrow e
\]
le foncteur à valeurs dans la catégorie ponctuelle. On veut que
\[
\left\{
\begin{array}{rcccl}
\mathrm{hot}(X) \times_{\D} M & \simeq & p_!(p^*(M)) & = & H^{\D}_{\bullet}(X, M)\\
\hm_{\D}(\mathrm{hot}(X), M) & \simeq & p_*(p^*(M)) & = & H^{\bullet}_{\D}(X, M)
\end{array}
\right.
\tag*{(4)}
\]
(isomorphismes canoniques).

\msp{page 3}
Ceci donnerait un principe de définition pour les opérations cherchées. Une difficulté vient du fait qu'en général, Diag n'est pas égal à (Cat), et que les types d'homotopie provenant de Diag ne donnent pas (même à isomorphisme près) \textit{tous} les types d'homotopie. Intuitivement, on ne trouve que des types d'homotopie `assez petits'. Mais on peut localiser Diag par rapport aux équivalences faibles (il vaut mieux dire `quasi-isomorphismes'), on trouve alors une catégorie $(\Hot)_0$, et un foncteur canonique
\[
\Hot_0 \longrightarrow \Hot,
\]
si on a de la chance, ce foncteur sera pleinement fidèle, et on essayera de décrire les foncteurs (1) ou (2) en y remplaçant Hot par la sous-catégorie pleine des `petits' types d'homotopie. On peut aussi oublier (provisoirement!) Hot, et travailler avec $\Hot_0$. C'est ce qu'on va faire pour le moment.

Pour pouvoir définir les foncteurs
\[
\left\{
\begin{array}{lccccccl}
\Hot_0 & \times & \mathcal{A} & \longrightarrow & \mathcal{A} & (\xi, F) & \longmapsto & \xi \times_{\D} F\\
\Hot_0^o & \times & \mathcal{A} & \longrightarrow & \mathcal{A} & (\xi, F) & \longmapsto & \hm_{\D}(\xi, F),
\end{array}
\right.
\tag*{(5)}
\]
je note que pour $F$ \texttt{[plutôt} $M$\texttt{]} fixé, $X$ variable dans Diag, les
\printedpage{3}
% head: A. Grothendieck : Dérivateurs, Ch. III, Édition des Universités de Montpellier II et Paris VII
% foot: Version 2 février 2002
\msp{page 4}
derniers membres dans (4)
\[
H^{\D}_{\bullet}(X, M), \qquad H^{\bullet}_{\D}(X, M)
\tag*{(6)}
\]
sont bien, l'un covariant, l'autre contravariant en $X$. Donc il y a , à isomorphisme unique près, une seule paire de bifoncteurs (5) satisfaisant (bifonctoriellement) les isomorphismes (4), si et seulement si les expressions (6), considérées comme fonctorielles en $X$, transforment quasi-isomorphismes en isomorphismes. Cela équivaut donc à la validité du % [sic] "il y a ," -> "il y a,", espace imprimée avant la virgule

\textbf{Théorème} (?) (\footnote{Nouvel axiome ?}). \textit{Soit $f : X' \longrightarrow X$ dans Diag. Si $f$ est un quasi-isomorphisme, c'est une $\D$-équivalence.}

Quand $X$ est la catégorie ponctuelle, ce n'est là autre que l'axiome Der 3 (p. 67). J'ignore si c'est une conséquence des axiomes Der 1 à Der 5 (d). Si ce n'est pas le cas, il faudrait l'ajouter aux axiomes, je crois. Je n'ai pas même prouvé qu'une équivalence de catégories dans Diag (au sens des catégories!) est une $\D$-équivalence!

Notons la \textit{formule d'adjonction}
\[
\mathrm{Hom}(\xi \times_{\D} M, N) \xrightarrow{\sim} \mathrm{Hom}(M, \hm_D(\xi, N)).
\tag*{(7)}
\] % [sic] \hom_D (D italique maigre) -> \hom_{\mathbf{D}}, le dérivateur est noté en gras partout ailleurs
Si $\xi$ provient de $X \in \mathrm{Diag}$, alors cela s'écrit
\[
\mathrm{Hom}(p_!p^*M, N) \xrightarrow{\sim} \mathrm{Hom}(M, p_*p^*N),
\tag*{(8)}
\]
\msp{page 5}
or par adjonction les deux membres sont canoniquement isomorphes à
\[
\mathrm{Hom}_{\D(X)}(p^*M, p^*N).
\]

\textbf{Exemple.} Prenant pour $\xi$ le type d'homotopie du cercle standard (représenté p.ex. par l'ensemble ordonné
\[
\begin{tikzcd}[row sep=small, column sep=small]
 & b & \\
0 \arrow[ur] \arrow[dr] & & 0' \arrow[ul] \arrow[dl] \ , \\
 & a &
\end{tikzcd}
\]
où on convient de prendre 0 comme origine, et d'orienter de 0 vers $a$), les deux foncteurs en question sont (sauf erreur) les foncteurs suspension $\Sigma$ et `cosuspension' ou foncteur `espace des lacets' $\Omega$. Sous réserve de vérifier que $\Omega F$ coïncide bien avec ce que j'ai décrit ailleurs, à propos des carrés substantiels cartésiens dans $\mathcal{A}$.

Dans $\Hot_0$ on a une opération produit, en fait ça doit être le produit catégorique dans $\Hot_0$. Mais sur les modèles il est donné par
\[
\mathrm{hot}_0(X) \times \mathrm{hot}_0(Y) \ \simeq \ \mathrm{hot}_0(X \times Y)
\tag*{(8)}
\] % [sic] (8) -> (9), le numéro (8) est déjà employé plus haut sur cette page

% ---- from Der-ChIII-567_p006-010.tex
\printedpage{4}
% head: A. Grothendieck : Dérivateurs, Ch. III, Édition des Universités de Montpellier II et Paris VII
% foot: Version 2 février 2002

\texttt{[deux fois} (8)\texttt{]}. Ceci posé, on a les isomorphismes canoniques, pour $\xi,\eta\in\Ob\Hot_0$,
\[
\left\{
\begin{aligned}
\xi\times_{\D}(\eta\times_{\D}M) &\;\simeq\; (\xi\times\eta)\times_{\D}M\\
\hm_{\D}(\xi,\hm_{\D}(\eta,M)) &\;\simeq\; \hm_{\D}(\xi\times\eta,M).
\end{aligned}
\right.
\tag*{(9)}
\]

Si $\xi$, $\eta$ sont représentés par $X$, $Y$, la première formule p.ex.
\msp{page 6}
s’établit en calculant le deuxième membre comme
\[
\varphi_!(M_{X\times Y})\simeq p_!(\;\underbrace{q'_!(M_{X\times Y})}_{\substack{\simeq\ q_!(M_Y)_X\\ \text{car $q$ est $\mathbf{D}$-lisse}}}\;)\simeq\xi\times_{\D}(\eta\times_{\D}M),
\]
O.K.
\[
\begin{tikzcd}
X \ar[d,"p"'] & X\times_e Y \ar[l,"q'"'] \ar[d,"p'"] \ar[dl,dotted,"\varphi"'] \\
e & Y \ar[l,"q"]
\end{tikzcd}
\]

J’ai aussi envie de prouver :

\textbf{Proposition 1.} \textit{Soit $\xi\in\Ob\Hot_0$. Alors $F\longmapsto\xi\times_{\D}F$, en tant que morphisme de dérivateurs $\D\longrightarrow\D$, est} exact à droite, \textit{et $F\longmapsto\hm_{\D}(\xi,F)$ est exact à gauche.}

On peut supposer que $\xi$ est défini par $X$ dans Diag. Soit $f:I\longrightarrow J$, soit $F$ dans $\D(I)$, il faut prouver
\[
f_*(\hm_{\D}(\xi,F))\;\simeq\;\hm_{\D}(\xi,f_*(F)).
\]
Considérons donc
\[
\begin{tikzcd}
X\times I \ar[r,"f'"] \ar[d,"{p=(x,i)\longmapsto i}"'] & X\times J \ar[d,"{q=(x,j)\longmapsto j}"] \\
I \ar[r,"f"] & J.
\end{tikzcd}
\]

On a
\[
\hm_{\D}(\xi,F)=p_*(p^*F),
\]
donc
\[
f_*(\hm_{\D}(\xi,F))=\underbrace{f_*p_*}_{(fp)_*}(p^*(F))=(qf')_*p^*(F)\simeq q_*(f'_*p^*(F)).
\]
Mais
\[
f'_*p^*(F)\;\simeq\;q^*f_*(F),
\]
formule de changement de base des $f_*$, pour le changement de base $q$, qui est lisse, car déduit de $X\longrightarrow e$, qui est lisse, par changement de base $J\longrightarrow e$. On trouve
\printedpage{5}
% head: A. Grothendieck : Dérivateurs, Ch. III, Édition des Universités de Montpellier II et Paris VII
% foot: Version 2 février 2002
\msp{page 7}
donc
\[
f_*(\hm_{\D}(\xi,F))\simeq q_*q^*(f_*(F))\defeq\hm_{\D}(\xi,f_*(F)),
\]
q.e.d. L’énoncé pour $F\longmapsto\xi\times_{\D}F$ est dual de celui qu’on vient de prouver.

Reprenons, pour $\xi\in(\mathrm{hot})_0$ \texttt{[}$=\Hot_0$\texttt{]}, la formule d’adjonction
\[
\Hom(\xi\times_{\D}M,N)\simeq\Hom(M,\hm_{\D}(\xi,N))\defeq\Hom_{\xi}(M,N).
\tag*{$(*)$}
\]

Si $\xi$ est défini par $X$ dans Diag, les deux membres de $(*)$ s’identifient à
\[
\Hom_{\D(X)}(p^*M,p^*N),
\]
où $p:X\longrightarrow e$ est le foncteur structural.

Fixons $M$, $N$, et considérons l’expression $(*)$ comme contrafoncteur
\[
(\Hot_0)^o\;\longrightarrow\;\mathrm{Ens}.
\]

Il semblerait que dans tous les cas connus, ce foncteur soit ‘représentable’ par un objet, sinon de $\Hot_0$ lui-même, du moins un objet dans Hot. Si on ne peut le définir comme objet $H$ de $\Hot_0$, le fait qu’on ait un isomorphisme canonique
\[
\Hom_{\xi}(M,N)\;\simeq\Hom_{\Hot}(\xi,H),
\tag*{(10)}
\]
fonctoriel en $\xi\in\Ob\Hot_0$, ne \texttt{[le]} définit pas forcément
\msp{page 8}
à isomorphisme unique près. En fait, quand
\[
\D=\D_{\mathcal{M},\Sigma}
\]
est défini par une catégorie de modèles $(\mathcal{M},\Sigma)$, alors il y a une définition naturelle d’une opération $\homext$ (Hom ‘externe’) à valeurs dans Hot, savoir un foncteur
\[
\left\{
\begin{aligned}
&\mathcal{A}^o\times\mathcal{A} &&\longrightarrow &&\mathrm{Hot}\\
&(M,N) &&\longmapsto &&\homext(M,N),
\end{aligned}
\right.
\tag*{(11)}
\]
où $\mathcal{A}=\mathcal{M}\Sigma^{-1}$ est la catégorie fondamentale de $\D=\D_{\mathcal{M},\Sigma}$. Il est vrai que je n’ai pas vérifié, si p. ex. $(\mathcal{M},\Sigma)$ est sous-jacente à une catégorie de modèles quillénienne, de sorte que (par \textsc{Thomason}) $\D$ est bien un dérivateur, que cette notion de $\homext$ satisfait bien à une formule de la forme (10). Ici je vais le poser [?], au besoin, ces données supplémentaires : donnée de (11), et d’un isomorphisme fonctoriel
\[
\underbrace{\Hom_{\xi}(M,N)}_{\begin{subarray}{l}=\ \Hom(\xi\times_{\D}M,N)\\ =\ \Hom(M,\hm_{\D}(\xi,N))\end{subarray}}
\quad\simeq\;\Hom_{\Hot}(\xi,\homext(M,N)).
\tag*{(12)}
\]

\printedpage{6}
% head: A. Grothendieck : Dérivateurs, Ch. III, Édition des Universités de Montpellier II et Paris VII
% foot: Version 2 février 2002
\msp{page 9}
Si maintenant $\eta$ est un deuxième élément de $\Hot_0$, on a envie d’écrire
\[
\left\{
\begin{aligned}
\homext(\eta\times_{\D}M,N) &\;\simeq\;\homext(M,\hm_{\D}(\eta,N))\\
&\;\simeq\;\hm_{\HOT}(\eta,\homext(M,N)),
\end{aligned}
\right.
\tag*{(13)}
\]
où dans le deuxième membre, $\hm_{\HOT}$ est l’opération (1) dans le cas où $\D=\mathrm{HOT}$, et qui est défini aussi en termes du produit cartésien $\xi\times\eta$ dans Hot, par
\[
\Hom_{\Hot}(\xi,\hm_{\HOT}(\eta,\zeta))\;\simeq\;\Hom_{\Hot}(\xi\times\eta,\zeta),
\]
en d’autres termes, pour $\eta$ fixé, le foncteur
\[
\zeta\;\longmapsto\;\hm_{\HOT}(\eta,\zeta)
\]
est défini comme le foncteur adjoint à droite du foncteur
\[
\xi\;\longmapsto\;\xi\times\eta.
\]

Pour vérifier heuristiquement la formule (13) (vérification valable ‘en forme’ si les trois objets envisagés sont en fait définis comme objets de $\Hot_0\ \ldots$), on prend $\Hom_{\Hot}(\xi,-)$ à valeurs dans les trois membres, pour $\xi$ donné dans $\Hot_0$. Pour le premier, on trouve
\[
\begin{aligned}
\Hom_{\xi}(\eta\times_{\D}M,N) &\;=\;\Hom(\xi\times_{\D}(\eta\times_{\D}M),N)\\
&\;=\;\Hom((\xi\times\eta)\times_{\D}M,N).
\end{aligned}
\]
Pour le deuxième
\[
\begin{aligned}
\Hom_{\xi}(M,\hm_{\D}(\eta,N)) &\;=\;\Hom_{\mathcal{A}}(\xi\times_{\D}M,\hm_{\D}(\eta,N))\\
&\;\simeq\;\Hom_{\mathcal{A}}(\eta\times_{\D}(\xi\times_{\D}M),N)\\
&\;\simeq\;\Hom_{\mathcal{A}}((\eta\times\xi)\times_{\D}M,N),
\end{aligned}
\]
\msp{page 10}
et enfin pour le troisième
\[
\begin{aligned}
\Hom_{\Hot}(\xi\times\eta,\homext(M,N)) &\;\simeq\;\Hom_{\xi\times\eta}(M,N)\\
&\;\simeq\;\Hom((\xi\times\eta)\times_{\D}M,N),
\end{aligned}
\]
donc toujours la même chose ! Pas de doute, ça doit marcher.

S’il faut se donner le $\homext$ (11) comme donnée supplémentaire, avec de plus un isomorphisme trifonctoriel (12), il faut aussi se donner les isomorphismes trifonctoriels (13), \textit{‘compatibles’} avec l’isomorphisme (12) au sens du calcul précédent.

\printedpage{7}
% head: A. Grothendieck : Dérivateurs, Ch. III, Édition des Universités de Montpellier II et Paris VII
% foot: Version 2 février 2002

Par exemple, prenant $\eta=$ type d’homotopie du cercle $S^1$, les formules (13) se spécialisent en
\[
\begin{aligned}
\homext(\Sigma M,N) &\;\simeq\;\homext(M,\Omega N)\\
&\;\simeq\;\Omega\homext(M,N)\\
&\;\defeq\;\hm_{\HOT}(S^1,\homext(M,N)).
\end{aligned}
\tag*{(14)}
\]

Mais cette formule dans Hot n’est sans doute pas tout à fait la bonne, car le bon formalisme de $\Omega$, $\Sigma$ se fait dans les catégories \textit{ponctuées,} et en particulier dans la catégorie homotopique ponctué $\Hot_{\bullet}$. % [sic] ponctué -> ponctuée, agreement with « catégorie »

Je vais expliciter un exemple, en prenant le dérivateur
\[
\mathrm{Dér}(A)
\]
‘dérivé’ d’une catégorie abélienne $A$. Donc $\D(e)=$ catégorie dérivée de $A$. Pour éviter des
\msp{page 11}
difficultés accessoires ici, je vais supposer que $A$ est la catégorie des $k$-modules, $k$ un anneau quelconque.

À tout $\xi\in\Ob\Hot$ est associé un ‘abélianisé’
\[
\xi_{\ab}\defeq H_{\bullet}(\xi)\in\underbrace{D(\mathrm{Ab})}_{\text{ou }\mathbf{D}(\mathbf{Z}\text{-Mod})}\;,\qquad\quad\text{catégorie dérivée des groupes abéliens,}
\]
objet qui n’a de cohomologie qu’en degré $\le 0$. Ceci dit, si $M$ est dans $D(k\text{-Mod})$, on aura
\[
\left\{
\begin{aligned}
\xi\times_{\D}M &\;\simeq\;\xi_{\ab}\Ltens_{\mathbf{Z}}M\\
\hm_{\D}(\xi,M) &\;\simeq\;\RHom_{\mathbf{Z}}(\xi_{\ab},M),
\end{aligned}
\right.
\tag*{(15)}
\]
formule qui est en cheville avec la formule de Künneth pour l’homologie
\[
(\xi\times\eta)_{\ab}\simeq\xi_{\ab}\Ltens_{\mathbf{Z}}\eta_{\ab},
\tag*{(16)}
\]
ce qui redonne un aspect familier aux formules (9), qui se réduisent à
\[
\left\{
\begin{aligned}
\xi_{\ab}\Ltens_{\mathbf{Z}}(\eta_{\ab}\Ltens_{\mathbf{Z}}M) &\;\simeq\;(\xi_{\ab}\Ltens_{\mathbf{Z}}\eta_{\ab})\Ltens_{\mathbf{Z}}M\\
\RHom_{\mathbf{Z}}(\xi_{\ab},\RHom_{\mathbf{Z}}(\eta_{\ab},M)) &\;\simeq\;\RHom_{\mathbf{Z}}(\xi_{\ab}\Ltens_{\mathbf{Z}}\eta_{\ab},M).\qquad\text{(\footnotemark[2])}
\end{aligned}
\right.
\]
\footnotetext[2]{Explicitier aussi la formule d’adjonction (7) (p. 4).} % [sic] Explicitier -> Expliciter, spelling (cf. « expliciter » in the text)

D’autre part, existe-t-il bien un $\homext(M,N)$, si $M$, $N$ dans $D(k\text{-Mod})$ ? On pense à $\RHom_k(M,N)$, mais pour en faire un type d’homotopie, on doit ‘tronquer en degrés $>0$’. On pourrait donc prendre
\[
\homext(M,N)\;\simeq\;\tau_{\le 0}(\RHom_k(M,N)).
\]

\printedpage{8}
% head: A. Grothendieck : Dérivateurs, Ch. III, Édition des Universités de Montpellier II et Paris VII
% foot: Version 2 février 2002
\msp{page 12}
La formule fondamentale (10) est-elle alors valable ? Le deuxième membre est en effet
\[
\Hom_{\Hot}(\xi,H)\;\simeq\;\Hom_{\mathrm{Hotab}}(\xi_{\ab},H)\qquad\quad\text{(\footnotemark[3])}
\]
\footnotetext[3]{? ? est-ce vrai ? ? ?}
car $H$ est un objet abélien de Hot (cf. le tapis sur l’abélianisation \ldots), à partir de là c’est gagné :
\[
\begin{aligned}
\Hom_{\underbrace{\scriptstyle\mathrm{Hotab}}_{\begin{subarray}{l}\approx\ D(\mathbf{Z}\text{-Mod})^{\le 0},\\ \text{sous-catégorie \textit{pleine}}\\ \text{de }D(\mathbf{Z}\text{-Mod})\end{subarray}}}(\xi_{\ab},\tau_{\le 0}\RHom_k(M,N)) &\;\simeq\;\Hom_{D(\mathbf{Z}\text{-Mod})}(\xi_{\ab},\RHom_k(M,N))\\
&\;\simeq\;\Hom_{D(k\text{-Mod})}(\xi_{\ab}\times M,N),
\end{aligned}
\]
O.K.

D’autre part, le bifoncteur
\[
(M,N)\;\longmapsto\;\homext(M,N)\;=\;\tau_{\le 0}\RHom_k(M,N)
\]
se factorise canoniquement par Hotab et a fortiori par $\Hot_{\bullet}$. Le ‘bon’ $\homext(M,N)$ serait ici plutôt celui dans Hotab, voire même dans la catégorie $D^{-}(\mathbf{Z}\text{-Mod})$, déduite de celle là ‘en rendant le foncteur translation inversible’. Il est évident d’ailleurs qu’ici le $\xi_{\ab}\times_{\D}M$, où $\xi=S^1$, n’est pas ‘la bonne’ suspension, qui doit être le décalage des degrés – $\xi_{\ab}$ ne devrait avoir de la cohomologie qu’en degré $-1$, et non aussi en degré $0$.

% ---- from Der-ChIV-567_p001-005.tex
\printedpage{}% no page number printed on this page
% head: (none)

\begin{center}
{\LARGE\textsf{LES DÉRIVATEURS}}

\bigskip

\textbf{ALEXANDRE GROTHENDIECK}

\bigskip\bigskip

\textbf{Chapitre IV}

\medskip

\textbf{Diagrammes substantiels}
\end{center}

\clearpage

\printedpage{}% no page number printed on this page
% head: (none)

\noindent Ce texte a été déchiffré et transcrit en \LaTeXe\ par M. Künzer. Il a été édité par M. Künzer, J. Malgoire, et G. Maltsiniotis. La transcription est aussi fidèle que possible au manuscrit. Pour les quelques corrections évidentes, ou rares commentaires des éditeurs, ainsi que pour la numérotation originale des pages du manuscrit, les caractères de machine à écrire \texttt{[typewriter]} entre crochets sont utilisés. Un point d'interrogation entre crochets signifie que l'on n'est pas sûr du mot qui précède.

\noindent Cette édition est provisoire. Les remarques, commentaires et corrections sont bienvenus. Envoyer un message à :

\noindent\textbf{[e-mail address removed]}

\hfill G. Maltsiniotis\hspace*{2cm}

\clearpage

\printedpage{1}
% head: A. Grothendieck : Dérivateurs, Ch. IV, Édition des Universités de Montpellier II et Paris VII
% foot: Version 2 février 2002

\noindent\texttt{[Il n'y a pas de pages 1-72]}

\msp{page 73}

\subsection*{Construction de diagrammes substantiels pour un dérivateur.}

\noindent Soit $\D$ un prédérivateur, $I$ dans Diag, considérons

\noindent $\varphi_I : \mathcal{D}(I) \longrightarrow \uHom(I^o,\cA)$,\ \ où $\cA=\D(e)$ est la catégorie fondamentale des coefficients. % [sic] \mathcal{D}(I) -> \mathbf{D}(I), elsewhere the prederivator is the bold D

Les objets de $\uHom(I^o,\cA)$ s'interprètent comme des \textit{diagrammes de type $I^o$} dans $\cA$. Ce langage est surtout commode si $I$ est ordonnée. Notons, pour $F\in\mathrm{Ob}\,\D(I)$, par $(F_i)_{i\in\mathrm{Ob}\,I}$ le $I^o$-diagramme associé dans $\cA$. Par l'axiome Dér 2, le foncteur
\[
\varphi : F \longmapsto (F_i)_{i\in\mathrm{Ob}\,I}
\]
est conservatif. Mais il n'a aucune envie d'être fidèle, encore moins pleinement fidèle. Quand on connaît le diagramme $(F_i)_i$, on est loin de connaître $F$, pas même à isomorphisme non unique près, je présume. Il n'est nullement clair que si $F$, $F'$ dans $\D(I)$ sont tels que les diagrammes associés soient isomorphes, ils soient eux-mêmes isomorphes.

Pour insister sur ce point, les objets de $\D(I)$ s'appelleront les `\textit{diagrammes substantiels} de type $I^o$ dans $\cA$' (ou $\D$-substantiel, s'il y a risque de confusion). Pour étayer l'intuition et simplifier les explications, un objet $F$ de $\D(I)$ ($\D$-coefficients sur $I$) sera représenté graphiquement par le diagramme `faible' ou `squelettique' correspondant -- on le regarde
\msp{page 74}
comme une sorte de (pâle) `squelette' du diagramme substantiel dont il provient. P. ex. si $I=\bDelta_0 : 0\longrightarrow 1$ \texttt{[plutôt $\bDelta_1$]}, on note un diagramme substantiel de type $(\bDelta_0)^o$ \texttt{[plutôt $(\bDelta_1)^o$]} par le même graphisme, que le diagramme ordinaire :
\[
F_0 \longleftarrow F_1.
\]
Si
\[
I = \begin{tikzcd}
a \arrow[r] & c \\
0 \arrow[u] \arrow[r] & b, \arrow[u]
\end{tikzcd}
\]
les diagrammes substantiels de type $I$ se notent
\[
\begin{tikzcd}
F_a \arrow[d] & F_c \arrow[l] \arrow[d] \\
F_0 & F_b. \arrow[l]
\end{tikzcd}
\]

\clearpage

\printedpage{2}
% head: A. Grothendieck : Dérivateurs, Ch. IV, Édition des Universités de Montpellier II et Paris VII
% foot: Version 2 février 2002

Il faudra bien sûr préciser chaque fois si ce graphisme désigne un diagramme squelettique, ou substantiel. Souvent on part d'un squelette, et on essaye de `l'étoffer' avec un peu de chair, pour en faire un diagramme substantiel -- si possible, de façon canonique. C'est là un point de philosophie dérivatique essentiel pour la pratique. On procédera souvent de proche en proche, en `substantifiant' des sous-diagrammes correspondants à une suite strictement croissante $I_0\subset I_1\subset\cdots\subset I_n=I$ de $I$, au mieux de la situation particulière.

\msp{page 75}

En voici un exemple typique, en supposant $\D$ un dérivateur \textit{ponctué} (les $\D(I)$ sont ponctuées, et les $f^*$ transforment objets nuls en objets nuls). Il s'agit de la définition du foncteur `cosuspension' ou `espace de lacets' dans $\cA$,
\[
\Omega : \cA \longrightarrow \cA.
\]
Il est entendu que pour $M$ dans $\cA$, $\Omega(M)$ doit s'insérer dans un carré commutatif
\[
\begin{tikzcd}
e \arrow[d] & \Omega M \arrow[l] \arrow[d] \\
M & e, \arrow[l]
\end{tikzcd}
\tag*{(C)}
\]
où $e$ désigne les éléments unités. Ce carré cependant devrait être substantiel, plus précisément, être un carré substantiel `$\D$-\textit{cartésien}'. Étant entendu qu'un objet de $\D(I)$ (carré substantiel dans $\cA$) est dit $\D$-cartésien, s'il est dans l'image essentielle de $\D(I_1)$ par $i_{1*}$, où
\[
i_1 : I_1 \hookrightarrow I,
\]
$I_1$ étant la sous-catégorie ouverte de $I$ définie par
\[
I_1 = \begin{tikzcd}
a \\
0 \arrow[u] \arrow[r] & b.
\end{tikzcd}
\]

\msp{page 76}

La restriction du carré substantiel cherché (C) à la sous-catégorie $I_1$ est donc un diagramme substantiel de squelette
\[
\begin{tikzcd}
e \arrow[d] \\
M & e, \arrow[l]
\end{tikzcd}
\tag*{(Sq)}
\]
\printedpage{3}
% head: A. Grothendieck : Dérivateurs, Ch. IV, Édition des Universités de Montpellier II et Paris VII
% foot: Version 2 février 2002
dont il se déduit (à isomorphisme canonique près) en appliquant $i_*$ à ce dernier, de façon à le compléter en $\begin{tikzcd}
e \arrow[d] & \Omega M \arrow[l, dotted] \arrow[d, dotted] \\
M & e \arrow[l]
\end{tikzcd}$. Il reste donc à décrire une substantiation canonique du squelette (Sq), en termes de $M$. Cela nous conduit à regarder la sous-catégorie ouverte
\[
I_0 = \{0\} :\quad
\begin{tikzcd}
{} \arrow[r, dotted] & {} \\
0 \arrow[u, dotted] \arrow[r, dotted] & {} \arrow[u, dotted]
\end{tikzcd}
\]
de $I_1$ (donc aussi de $I$) réduite à l'objet initial, et l'inclusion
\[
I_0 \overset{i_0}{\hookrightarrow} I_1,
\]
et de substantifier le squelette (Sq) par $(i_0)_!(M)$.

On peut dire a posteriori que, si on connaissait le carré substantiel, il n'y a, à isomorphisme unique près, qu'une seule façon de l'obtenir comme un $i_{1*}i_{0!}(?)$, avec ? dans $\D(I_0)\simeq\cA$, car $i_{0!}$ et $i_{1*}$ étant pleinement fidèles, il en est de même de leur composé.

\msp{page 77}

En fait, on voit que le foncteur
\[
i_{1*}i_{0!} :\ \cA\simeq\D(I_0) \longrightarrow \D(I)
\]
induit une équivalence de $\cA$ avec la sous-catégorie pleine de $\D(I)$ formée des carrés $\D$-cartésiens dont les fibres en $a$ et $b$ sont nulles. (Le foncteur quasi-inverse évident est $i^*$, où $i : U_0 \hookrightarrow I$ \texttt{[plutôt $I_0 \hookrightarrow I$]}.) Cela signifie donc que dans ce cas, la donnée du squelette de nature très particulière ($F_a\simeq F_b\simeq 0_{\cA}$), plus la contrainte supplémentaire $\D$-cartésienne, sur le carré cherché, détermine celui-ci à isomorphisme unique près en termes de son squelette : il y a une substantiation non seulement `canonique' de celui-ci, mais même unique à isomorphisme unique près, à cause de la contrainte $\D$-cartésienne.

Je voudrais commencer par regarder le cas plus simple encore où $I=\bDelta_1=(0\longrightarrow 1)$, donc les squelettes envisagés sont simplement les flèches
\[
F_1 \longrightarrow F_0 \qquad \text{dans } \cA.
\]
La philosophie nous oblige, pour les flèches les plus fondamentales dans la théorie des dérivateurs $\D$ (flèches dans $\cA$, ou par extension dans une $\D(I)\simeq\cA_{\D_I}$, catégorie fondamentale du dérivateur induit $\D_I$) de se demander s'il n'y a pas

% ---- from Der-ChIV-567_p006-009.tex
\printedpage{4}
% head: A. Grothendieck : Dérivateurs, Ch. IV, Édition des Universités de Montpellier II et Paris VII
% foot: Version 2 février 2002
\msp{page 78}
une façon ‘canonique’, ‘naturelle’, ‘voulue par Dieu’, de le ‘remplumer’ en un diagramme substantiel. Les flèches les plus évidentes qui se présentent sont relatives à un foncteur
\[ f : I \longrightarrow J \]
dans Diag, et aux deux morphismes d'adjonction pour le couple $(f^*, f_*)$, savoir
\[
\begin{aligned}
G &\longrightarrow f_*f^*G && \text{dans } D(J)\\
f^*f_*F &\longrightarrow F && \text{dans } D(I). \quad (\text{\footnotemark})
\end{aligned}
\]
\footnotetext{Bien sûr, problème dual pour $(f_!, f^*)$.}

Je vais essayer de remplumer ces deux flèches, et vais généraliser en envisageant la situation ci-contre \texttt{[ci-dessous]}
\[
\begin{tikzcd}
I \arrow[d, "f"'] & F \\
J & G, \arrow[u, "u"']
\end{tikzcd}
\]
où
\[ \Hom^f(G,F) \overset{\text{déf}}{=} \Hom(f^*G,F) \overset{\text{adj.}}{\simeq} \Hom(G,f_*F). \]
Donc le même $u$ s'interprète par \textit{deux} flèches squelettiques
\[
\left\{
\begin{aligned}
f^*(G) &\longrightarrow F && \text{dans } \mathbf{D}(I)\\
G &\longrightarrow f_*(F) && \text{dans } \mathbf{D}(J).
\end{aligned}
\right.
\tag*{(1)}
\]
Les remplumer, c'est les interpréter (si faire se peut) en termes de deux objets (si possible canoniques) dans $\mathbf{D}(I \times \bDelta_1)$ et $\mathbf{D}(J \times \bDelta_1)$ respectivement.

Je vais me guider sur le cas particulier d'un dérivateur trivial
\[ \mathbf{D} = \mathbf{D}_{\mathcal{M}}, \]
où $\mathcal{M}$ est une catégorie de modèles ‘nue’ $(\Sigma = \emptyset)$.

\msp{page 79}
Pour $f$ fixé, on va construire une catégorie auxiliaire $K(f)$, de telle façon que pour toute catégorie $\mathcal{M}$, les systèmes $(G,F,u)$ d'un $\mathcal{M}$-préfaisceau $G$ sur $J$, d'un autre $F$ sur $I$, et d'un $f$-homomorphisme $G \longrightarrow F$, s'identifient (fonctoriellement) aux $\mathcal{M}$-préfaisceaux sur
\printedpage{5}
% head: A. Grothendieck : Dérivateurs, Ch. IV, Édition des Universités de Montpellier II et Paris VII
% foot: Version 2 février 2002
$K(f)$. Ceci dit, on va essayer d'interpréter les flèches (1) squelettiques, mais peu importe ici -- tous les $\mathbf{D}$-squelettes sont maintenant remplumés) en termes de la géométrie de $K(f)$. % [sic] remplumés) -> remplumés, closing parenthesis printed without any opening parenthesis

\[
\left\{
\begin{aligned}
&\Ob K(f) = \Ob K(I) \amalg \Ob K(J) \qquad \text{\texttt{[plutôt} $\Ob I \amalg \Ob J$\texttt{ ]}}\\
&\text{$I$, $J$ s'identifiant à des sous-catégories pleines de $K(f)$,}\\
&\text{il n'y a pas des flèches d'un $j \in J$ dans un $i \in I$.}\\
&\text{(Donc $I$ est ouverte dans $K(f)$, $J$ son complémentaire fermé.)}\\
&\text{Pour $i \in \Ob I$, $j \in \Ob J$, on pose $\Hom_{K(f)}(i,j) = \Hom_J(f(i),j)$.}
\end{aligned}
\right.
\]

La composition des flèches dans $K(f)$ est évidente. On va définir une rétraction canonique
\[ \bar{f} : K(f) \longrightarrow J \]
de $K(f)$ sur sa sous-catégorie fermée $J$. Donc $\bar{f}$ est connu sur $\Ob J$, et sur $\Ob I$ ce sera $f$ :
\[
\left\{
\begin{aligned}
\bar{f}(j) &= j && \text{si } j \in \Ob J\\
\bar{f}(i) &= f(i) && \text{si } i \in \Ob I.
\end{aligned}
\right.
\]
Il reste à définir, pour $x, y \in \Ob K(f)$,
\msp{page 80}
\[ \Hom_{K(f)}(x,y) \longrightarrow \Hom_J(\bar{f}(x), \bar{f}(y)). \]
Si $x, y \in \Ob J$, c'est l'identité, si $x, y \in I$, c'est $f$. Si $x \in I$, $y \in J$, on doit définir une flèche
\[ \Hom(f(x),y) \longrightarrow \Hom(f(x),y), \]
et on prend encore l'identité. On a
\[
\begin{tikzcd}
J \arrow[r, hook, "j"] \arrow[dr, equal] & K(f) \arrow[d, "\bar{f}"] & I \arrow[l, hook', "i"'] \\
& J &
\end{tikzcd}
\]
\[
\left\{
\begin{aligned}
\bar{f} \circ j &= \id_J\\
j \circ \bar{f} &= \left\{
\begin{aligned}
&j && \text{sur } J\\
&jf && \text{sur } I,
\end{aligned}
\right.
\end{aligned}
\right.
\qquad \text{i.e.} \quad
\boxed{
\begin{aligned}
j\bar{f}j &= j && \text{sur } J\\
j\bar{f}i &= jf && \text{sur } I
\end{aligned}
}\ ,
\]
et on a \texttt{[}$K = K(f)$\texttt{]}
\[
\begin{aligned}
\alpha :\ \id_K &\longrightarrow j\bar{f} && \text{i.e. homomorphisme fonctoriel}\\
x &\xrightarrow{\alpha(x)} j(\bar{f}(x)), && x \in \Ob K,
\end{aligned}
\]
\printedpage{6}
% head: A. Grothendieck : Dérivateurs, Ch. IV, Édition des Universités de Montpellier II et Paris VII
% foot: Version 2 février 2002
défini ainsi :
\[
\left\{
\begin{aligned}
&\text{Si } x \in J,\ \alpha(x) = \id_x\ :\ x \xrightarrow{\sim} j(\bar{f}(x)) = j(x) = x\\
&\text{Si } x \in I,\ \alpha(x) : \underbrace{x \longrightarrow f(x)}_{\text{dans } K(f)},\ \text{ou } \underbrace{f(x) \longrightarrow f(x)}_{\text{dans } J},\ \alpha(x) = \id_{f(x)}.
\end{aligned}
\right.
\]
On doit vérifier la fonctorialité pour $x$ variable dans $K$ \dots

Soit $\Phi$ un $\mathcal{M}$-préfaisceau sur $K$. Il définit
\[
\begin{aligned}
F &= i^*(\Phi) && \in \Ob \mathbf{D}_{\mathcal{M}}(I)\\
G &= j^*(\Phi) && \in \Ob \mathbf{D}_{\mathcal{M}}(J).
\end{aligned}
\]
D'autre part $f^*(G)$ s'interprète comme $(j\bar{f})^*(\Phi)|I$, car
\[
\begin{aligned}
f^*(G)(i) &= G(\ \underbrace{f(i)}_{= \bar{f}(i) \,\in\, J}\ )\\
&= j^*(\Phi)(\bar{f}(i))\\
&= \Phi(j\bar{f}(i))\\
&= (j\bar{f})^*(\Phi)(i).
\end{aligned}
\]

\msp{page 81}
Donc
\[ f^*(G) \simeq (j\bar{f})^*(\Phi)|I. \]
Or la flèche
\[ \alpha : \id_K \longrightarrow j\bar{f} \]
donne
\[ \underbrace{\id_K(\Phi)}_{= \Phi} \xleftarrow{\Phi(\alpha)} (j\bar{f})^*(\Phi), \qquad \text{fonctoriel en } \Phi, \]
d'où, en prenant les restrictions à $I$,
\[ i^*(\Phi) \overset{\text{déf}}{=} F \xleftarrow{\ u\ } f^*G. \]

En fait, sans tellement se tortiller, on devrait pouvoir montrer que pour toute catégorie $\mathcal{M}$, la catégorie $\uHom(K^o, \mathcal{M})$ des $\mathcal{M}$-préfaisceaux sur $K$ est isomorphe à celle des triples $(F,G,u)$, avec $F \in \Ob \uHom(I^o, \mathcal{M})$, $G \in \uHom(J^o, \mathcal{M})$, et $u : f^*(G) \longrightarrow F$. Ainsi, $K$ apparaît comme solution d'un problème universel
\[ \uHom(K^o, \mathcal{M}) \longrightarrow \{F,\ G,\ u \text{ comme dessus}\}(\mathcal{M}). \]

Mais pour y préciser la flèche $\longrightarrow$ de façon conceptuelle simple, il est bon de se donner $K$ avec les structures supplémentaires \begin{tikzcd}[cramped, column sep=normal] I \arrow[r, hook, shift left, "j"] & K \arrow[l, shift left, "\bar{f}"] & I \arrow[l, hook', "i"'] \end{tikzcd}, $\bar{f}j = \id_J$, $\id_K \xrightarrow{\alpha} j\bar{f}$. Et ce système % [sic] I (first object, source of j and target of f-bar) -> J, since j : J -> K(f) and f-bar : K(f) -> J
\printedpage{7}
% head: A. Grothendieck : Dérivateurs, Ch. IV, Édition des Universités de Montpellier II et Paris VII
% foot: Version 2 février 2002
\msp{page 82}
$\{K,\ i,\ j,\ \bar{f},\ \alpha\}$ est (sauf erreur) défini à isomorphisme canonique près par la condition de représenter le foncteur $\mathcal{M} \longmapsto \{F,\ G,\ u\}^o(\mathcal{M})$ \dots

En particulier, le système défini par $G \in \Ob \mathbf{D}_{\mathcal{M}}(J)$
\[ (f^*(G),\ G,\ \id : f^*(G) \longrightarrow f^*(G)) \]
est défini sauf erreur par $\bar{f}^*(G)$ sur $K(f)$. On se propose [?] de trouver une interprétation similaire pour
\[ (F,\ f_*(F),\ f^*f_*F \longrightarrow F) \]
$(F \in \Ob \mathbf{D}(I))$. Le candidat naturel serait ici $\Phi = i_*(F)$. Admettons ce point (i.e. surtout la vérification du fait que $j^*(i_*(F)) \simeq f_*(F)$). Tout ceci montre que si $\mathbf{D}$ est à nouveau un dérivateur quelconque, on peut, par les formules précédentes, trouver un objet canonique $\Phi$ de $\mathbf{D}(K)$, en termes soit d'un $\mathbf{D}$-coefficient $F$ sur $I$ (en prenant $\Phi = i_*(F)$), soit d'un $\mathbf{D}$-coefficient $G$ sur $J$ (en prenant $\Phi = \bar{f}^*(G)$). D'autre part, on peut espérer définir canoniquement
\[
\begin{aligned}
\mathbf{D}(K) &\longrightarrow \mathbf{D}(I \times \bDelta_1)\\
\mathbf{D}(K) &\longrightarrow \mathbf{D}(J \times \bDelta_1)
\end{aligned}
\]
dont les composés avec les foncteurs ‘squelettes’
\[
\begin{aligned}
\mathbf{D}(I \times \bDelta_1) &\longrightarrow \uHom(\bDelta_1^o, \mathbf{D}(I))\\
\mathbf{D}(J \times \bDelta_1) &\longrightarrow \uHom(\bDelta_1^o, \mathbf{D}(J))
\end{aligned}
\]
\msp{page 83}
redonnent (pour $\Phi \in \Ob \mathbf{D}(K)$) les flèches squelettiques
\[
\left\{
\begin{aligned}
f^*(G) &\xrightarrow{\ u\ } F\\
G &\xrightarrow{\ v\ } f_*F
\end{aligned}
\right.
\]
associées à $(F,G,u)$. De façon précise, $f^*(G) \longrightarrow F$ se définit en termes de $\mathbf{D}$, et $\Phi \in \Ob \mathbf{D}(K)$, comme ci-dessus, et $v$ s'en déduit par adjonction.

L'idée naturelle ici, serait de ramener le problème de l'étoffement des flèches d'adjonction sqelettiques, au cas des immersions ouvertes ou fermées. Notons que les sous-catégories ouvertes de $K$ correspondent justement aux foncteurs $K \longrightarrow \bDelta_1$, en prenant les images inverses de l'ouvert générique. Il y a donc un unique % [sic] sqelettiques -> squelettiques, missing u (spelled squelettiques elsewhere in the text)
\[ \varphi :\ K(f) \longrightarrow \bDelta_1 \]
avec $I = \varphi^{-1}(\{0\})$. Donc \texttt{[fin du chapitre]}

% ---- from Der-ChV-234_p001-005.tex
\printedpage{}
% head: (none); no page number printed on this page

\begin{center}
\vspace*{6cm}

{\LARGE\textsf{LES DÉRIVATEURS}}

\vspace{1.5cm}

{\large\textbf{ALEXANDRE GROTHENDIECK}}

\vspace{2.5cm}

{\large\textbf{Chapitre V}}

\vspace{0.4cm}

{\large\textbf{Catégories de chemins et localisation}}
\end{center}

\clearpage

\printedpage{}
% head: (none); no page number printed on this page

\vspace*{6cm}

Ce texte a été déchiffré et transcrit en \LaTeXe{} par M. Künzer. Il a été édité par M. Künzer, J. Malgoire, et G. Maltsiniotis. La transcription est aussi fidèle que possible au manuscrit. Pour les quelques corrections évidentes, ou rares commentaires des éditeurs, ainsi que pour la numérotation originale des pages du manuscrit, les caractères de machine à écrire \texttt{[typewriter]} entre crochets sont utilisés. Un point d'interrogation entre crochets signifie que l'on n'est pas sûr du mot qui précède.

Cette édition est provisoire. Les remarques, commentaires et corrections sont bienvenus. Envoyer un message à :

\textbf{[e-mail address removed]}\\
\hspace*{\fill}G. Maltsiniotis\hspace*{2cm}

\clearpage

\printedpage{1}
% head: \textbf{A. Grothendieck :} \textit{Dérivateurs,} Ch. V, Édition des Universités de Montpellier II et Paris VII
% foot: Version 9 septembre 2001

\msp{page 1}

\section*{1\quad La catégorie des chemins}

Soit $\Phi$ la catégorie suivante :
\[
\left\{
\begin{array}{rcll}
\Ob\Phi & = & \mathbf{Z} & \\
\Fl\Phi & = & \Ob\Phi\sqcup\mathbf{Z} & \qquad (\Ob\Phi\leftrightarrow\text{identités})
\end{array}
\right.
\]
On note $\{\alpha\}$, ou simplement $\alpha$, l'objet qui correspond à $\alpha\in\mathbf{Z}$. On note $u_\alpha$ la flèche qui correspond à l'entier $\alpha\in\mathbf{Z}$.
\[
\begin{array}{rcl@{\qquad}l@{\quad}l@{\quad}l}
s(u_\alpha) & = & \alpha & \text{si } \alpha \text{ pair}, & \alpha+1 & \text{si } \alpha \text{ impair}\\[1ex]
b(u_\alpha) & = & \alpha+1 & \text{si } \alpha \text{ pair}, & \alpha-1 & \text{si } \alpha \text{ impair}
\end{array}
\]
\[
\texttt{[plutôt}\qquad b(u_\alpha) \;=\; \alpha+1 \quad \texttt{si}\ \alpha\ \texttt{pair,}\quad \alpha \quad \texttt{si}\ \alpha\ \texttt{impair]}
\]

Il n'y pas à définir de compositions, car si deux flèches sont composables, elles sont l'identité. En fait, $\Phi$ est une catégorie \textit{ordonnée}, $\{\alpha\}<\{\beta\}$ si et seulement si $\alpha$ pair, $\beta=\alpha+1$ \texttt{[plutôt} $\beta=\alpha+1$ \texttt{ou} $\beta=\alpha-1$\texttt{]}. % [sic] Il n'y pas -> Il n'y a pas, verb "a" missing

On note, pour $n\in\mathbf{N}$,
\[
\Phi_n\subseteq\Phi,\quad \Ob\Phi_n=\{\alpha\in\Ob\Phi\mid -n\le\alpha\le+n\},
\]
la sous-catégorie pleine de $\Phi$ dont les objets sont les $\{\alpha\}$ avec $-n\le\alpha\le+n$. Ainsi
\[
\Phi_1=\{-1\longleftarrow 0\longrightarrow 1\},
\]
la duale est
\[
\Psi_1=\Phi_1^{o}=\{-1\longrightarrow 0\longleftarrow 1\}.
\]
On peut considérer $\Phi$ comme somme amalgamée d'une suite ordonnée infinie (de deux côtés) de copies de $\Phi_1$, ou aussi d'une suite similaire de copies de $\Psi_1$.
\[
\begin{tikzcd}[column sep=small, execute at end picture={
\draw[decorate,decoration={brace,amplitude=4pt}] ([yshift=10pt]\tikzcdmatrixname-1-2.center) -- ([yshift=10pt]\tikzcdmatrixname-1-4.center) node[midway,above=4pt]{$\scriptstyle\Phi_1$};
\draw[decorate,decoration={brace,amplitude=4pt}] ([yshift=10pt]\tikzcdmatrixname-1-4.center) -- ([yshift=10pt]\tikzcdmatrixname-1-6.center) node[midway,above=4pt]{$\scriptstyle\Phi_1$};
\draw[decorate,decoration={brace,amplitude=4pt}] ([yshift=10pt]\tikzcdmatrixname-1-6.center) -- ([yshift=10pt]\tikzcdmatrixname-1-8.center) node[midway,above=4pt]{$\scriptstyle\Phi_1$};
\draw[decorate,decoration={brace,amplitude=4pt}] ([yshift=10pt]\tikzcdmatrixname-1-8.center) -- ([yshift=10pt]\tikzcdmatrixname-1-10.center) node[midway,above=4pt]{$\scriptstyle\Phi_1$};
\draw[decorate,decoration={brace,amplitude=4pt}] ([yshift=10pt]\tikzcdmatrixname-1-10.center) -- ([yshift=10pt]\tikzcdmatrixname-1-12.center) node[midway,above=4pt]{$\scriptstyle\Phi_1$};
\draw[decorate,decoration={brace,mirror,amplitude=4pt}] ([yshift=-10pt]\tikzcdmatrixname-1-3.center) -- ([yshift=-10pt]\tikzcdmatrixname-1-5.center) node[midway,below=4pt]{$\scriptstyle\Psi_1$};
\draw[decorate,decoration={brace,mirror,amplitude=4pt}] ([yshift=-10pt]\tikzcdmatrixname-1-5.center) -- ([yshift=-10pt]\tikzcdmatrixname-1-7.center) node[midway,below=4pt]{$\scriptstyle\Psi_1$};
\draw[decorate,decoration={brace,mirror,amplitude=4pt}] ([yshift=-10pt]\tikzcdmatrixname-1-7.center) -- ([yshift=-10pt]\tikzcdmatrixname-1-9.center) node[midway,below=4pt]{$\scriptstyle\Psi_1$};
\draw[decorate,decoration={brace,mirror,amplitude=4pt}] ([yshift=-10pt]\tikzcdmatrixname-1-9.center) -- ([yshift=-10pt]\tikzcdmatrixname-1-11.center) node[midway,below=4pt]{$\scriptstyle\Psi_1$};
}]
\cdots & {} & -4 \arrow[l] \arrow[r] & -3 & -2 \arrow[l] \arrow[r] & -1 & 0 \arrow[l] \arrow[r] & 1 & 2 \arrow[l] \arrow[r] & 3 & 4 \arrow[l] \arrow[r] & {} & \cdots
\end{tikzcd}
\]

Les $\Phi_n$, pour $n$ impair, sont de même somme de copies de $\Phi_1$ mises bout à bout, et les $\Psi_n$, pour $n$ pair $>0$, somme de copies de $\Psi_1$.

\clearpage

\printedpage{2}
% head: \textbf{A. Grothendieck :} \textit{Dérivateurs,} Ch. V, Édition des Universités de Montpellier II et Paris VII
% foot: Version 9 septembre 2001

Les $\Phi_n$ forment un système inductif avec foncteurs de transition pleinement fidèles, dont $\Phi$ est la limite inductive
\[
\Phi=\varinjlim_n\Phi_n.
\]
Mais pour la suite on va faire des $\Phi_n$ un système \textit{projectif}, en définissant pour tout $n\le m$ une \textit{rétraction} $\varphi_{n,m}$ de $\Phi_m$ sur $\Phi_n$, par

\msp{page 2}
\[
\left\{
\begin{array}{l}
\varphi_{n,m}:\Phi_m\longrightarrow\Phi_n\subset\Phi_m\\[1ex]
\varphi_{n,m}(\{\alpha_i\})=
\left\{
\begin{array}{rcll@{\quad}l@{\ \ }r}
\alpha_i & = & i & \in\Ob\Phi_n, & \text{i.e.} & -n\le i\le n\\[1ex]
\alpha_i & = & n & & \text{si} & n\le i\le m\\[1ex]
\alpha_i & = & -n & & \text{si} & -m\le i\le -n
\end{array}
\right.
\end{array}
\right.
\]

\texttt{[plutôt}
\[
\begin{array}{rcl}
\varphi_{n,m}:\Phi_m & \longrightarrow & \Phi_n\subset\Phi_m\\[1ex]
\varphi_{n,m}(\{i\}) & = &
\left\{
\begin{array}{l@{\qquad}l@{\quad}l@{\ \ }r}
\{i\} & \texttt{si}\ i\in\Ob\Phi_n, & \texttt{i.e} & -n\le i\le n\\[1ex]
\{n\} & & \texttt{si} & n\le i\le m\\[1ex]
\{-n\} & & \texttt{si} & -m\le i\le -n\ \texttt{]}
\end{array}
\right.
\end{array}
\]
On note le système projectif d'ensembles ordonnés ainsi obtenu par
\[
\bPhi=(\Phi_n)_{n\in\mathbf{N}}\in\mathrm{Pro}(\mathrm{Cat}).
\]
Soit $X$ une catégorie, alors pour $n\in\mathbf{N}$
\[
\uCh_n(X)\overset{\text{déf}}{=}\uHom(\Phi_n,X)
\]
est la catégorie des diagrammes de $X$ de la forme
\[
A_{-n}\cdots\longleftarrow A_{-2}\xrightarrow{u_{-2}}A_{-1}\xleftarrow{u_{-1}}A_0\xrightarrow{u_0}A_1\xleftarrow{u_1}A_2\xrightarrow{u_2}A_3\longleftarrow\cdots A_n
\tag*{$(*)$}
\]
avec $2n$ flèches bout à bout $u_i$ $(-n\le i\le n)$ \texttt{[plutôt (}$-n\le i\le n-1$\texttt{)]}, deux flèches consécutives allant toujours en sens opposé,
\[
\text{soit}\ \ \cdot\longrightarrow\cdot\longleftarrow\cdot,\ \ \text{soit}\ \ \cdot\longleftarrow\cdot\longrightarrow\cdot\,.
\]
Le foncteur canonique déduit de $\varphi_{n,m}:\Phi_m\longrightarrow\Phi_n$
\[
\varphi_X^{m,n}:\uCh_n(X)\longrightarrow\uCh_m(Y),\qquad n\le m,
\]
\printedpage{3}
% head: \textbf{A. Grothendieck :} \textit{Dérivateurs,} Ch. V, Édition des Universités de Montpellier II et Paris VII
% foot: Version 9 septembre 2001
\texttt{[plutôt} $\uCh_m(X)$\texttt{]} est le `foncteur de dégénérescence', consistant à compléter un diagramme $(*)$, par $m-n$ termes $A_{-n}$ avec des identités à gauche, et par $m-n$ termes \texttt{[}$A_n$\texttt{]} avec des identités à droite, donc pour $n$ pair on trouve
\[
\underbrace{A_{-n}\cdots\xrightarrow{\text{id}}A_{-n}}_{m-n\text{ termes}}\xleftarrow{\text{id}}A_{-n}\longrightarrow A_{-(n-1)}\cdots\cdots A_{n-1}\longleftarrow A_n\xrightarrow{\text{id}}\underbrace{A_n\xleftarrow{\text{id}}A_n\longrightarrow\cdots A_n}_{m-n\text{ termes}}.
\]
C'est un foncteur pleinement fidèle (ce qui correspond au fait que $\varphi_{n,m}$ est surjectif sur les objets et sur les flèches. % [sic] flèches. -> flèches)., opening parenthesis never closed

\msp{page 3}
On pose
\[
\underbrace{\uCh_\infty(X)}_{\text{ou simplement }\uCh(X)}\qquad=\uHom(\bPhi,X)\overset{\text{déf}}{=}\varinjlim_n\uCh_n(X)
\]
(\footnote{L'ensemble des objets de cette catégorie se notera $\mathrm{Ch}_\infty(X)$ ou $\mathrm{Ch}(X)$, de même on définit les $\mathrm{Ch}_n(X)$.}) (\footnote{dans cette formule: $\bPhi$ procatégorie}). On a une inclusion pleinement fidèle
\[
\uCh_\infty(X)\lhook\joinrel\longrightarrow\uHom(\Phi,X),
\]
où la deuxième catégorie est celle des diagrammes infinis (dans les deux directions)
\[
\cdots\longleftarrow A_{-2}\xrightarrow{u_{-2}}A_{-1}\xleftarrow{u_{-1}}A_0\xrightarrow{u_0}A_1\xleftarrow{u_1}A_2\xrightarrow{u_2}\cdots
\]
et où $\uCh_\infty(X)$ s'identifie à la sous-catégorie pleine formées des objets pour lesquels les flèches $u_i$ sont des identités, pour $|i|$ assez grand. % [sic] formées -> formée, agreement with "la sous-catégorie"

On appellera $\uCh_\infty(X)$ la \textit{catégorie des chemins} dans $X$, ses objets seront les chemins dans $X$. Ceux qui proviennent de $\uCh_n(X)$ seront appelés \textit{chemins standards de longueur} $2n$. (Le terme `standard' s'applique ici à nos conventions de notations. Ainsi nous ne considérons que des chemins de longueur paire, et pour une longueur déterminée, nous choisissons l'une parmi les deux espèces raisonnables, par exemple pour les chemins de longueur 2 nous avons pris seulement
\[
A_{-1}\longleftarrow A_0\longrightarrow A_1
\]
et non pas
\[
A_{-1}\longrightarrow A_0\longleftarrow A_1.)
\]

\msp{page 4}
Je me rends compte que mes notations $u_i$ ne sont pas commodes pour la dualité, je vais les changer, en $u_i$ pour les flèches entre objets $\{i\}$ avec $i\ge0$, et $v_i$ la flèche symétrique de $u_i$ par rapport à l'origine. (Le groupe des automorphismes des $\Phi_n$ est réduit au groupe à deux éléments, le seul automorphisme non trivial étant la symétrie $\{i\}\longrightarrow\{-i\}$. Donc nos chemins seront notés désormais % [sic] {-i}. -> {-i}). , opening parenthesis "(Le groupe" not closed on this page
\[
A_{-2}\xrightarrow{v_1}A_{-1}\xleftarrow{v_0}A_0\xrightarrow{u_0}A_1\xleftarrow{u_1}A_2\longrightarrow\cdots.
\tag*{$(*)$}
\]

% ---- from Der-ChV-234_p006-010.tex
\printedpage{4}
% head: A. Grothendieck : Dérivateurs, Ch. V, Édition des Universités de Montpellier II et Paris VII
% foot: Version 9 septembre 2001

(Ce n'est pas non plus enthousiasmant \dots)

Soit
\[
\varepsilon=\{-1,+1\}
\]
l'ensemble ordonné \textit{discret,} réduit aux deux éléments $-1$, $+1$, et définissons des foncteurs d'inclusion (pleinement fidèles)
\[
\begin{aligned}
\alpha_n:\varepsilon\longrightarrow\Phi_n\quad & \alpha_n(-1)=-n\\
& \alpha_n(1)=n,
\end{aligned}
\]
ces foncteurs sont compatible avec les $\varphi_{n,m}$, d'où % [sic] compatible -> compatibles, accord avec « ces foncteurs »
\[
\alpha_\infty:\varepsilon\longrightarrow\boldsymbol{\Phi}\in\mathrm{ProCat},
\]
d'où pour tout $X$ un foncteur
\[
\kappa:\uHom(\boldsymbol{\Phi},X)\longrightarrow\uHom(\varepsilon,X)\simeq X\times X
\]
i.e.

\msp{page 5}
\[
\begin{aligned}
\uCh(X) &\longrightarrow X\times X\\
c &\longrightarrow (\sigma(x),\beta(x)).
\end{aligned}
\]

\texttt{[plutôt} $(\sigma(c),\beta(c))$\texttt{]} Pour tout chemin $c$ dans $X$ i.e. un diagramme $(*)$ (satisfaisant à la condition de dégénérescence à l'infini), on a une source ou origine du chemin, et un but ou une \textit{extrémité}
\[
\left\{
\begin{aligned}
\sigma(x) &= A_i\quad\text{pour } i\leq -\mathrm{long}(x)\\
\beta(x) &= A_i\quad\text{pour } i\geq \mathrm{long}(x)
\end{aligned}
\right.
\]
\texttt{[remplacer} $x$ \texttt{par} $c$\texttt{; pour long}$(c)$ \texttt{défini convenablement]}. La fibre du foncteur précédent en $x$, $y$ sera notée
\[
\uCh(X;x,y)\qquad\text{ou simplement }\uCh(x,y),
\]
\textit{catégorie des chemins de $x$} à \textit{$y$.}

\printedpage{5}
% head: A. Grothendieck : Dérivateurs, Ch. V, Édition des Universités de Montpellier II et Paris VII
% foot: Version 9 septembre 2001

\section*{2\quad Calcul d'une catégorie de fractions $X\Sigma^{-1}$}

Avant de continuer le formalisme des chemins (composition etc.), je vais tester l'utilité de la notion de catégorie des chemins, pour le `calcul' des $X\Sigma^{-1}$, où
\[
\Sigma\subset\mathrm{Fl}(X),\qquad\Sigma\text{ contient les identités}
\]
est une partie autrement \textit{quelconque} de $\mathrm{Fl}(X)$. On sait que $\mathrm{Ob}\,X\Sigma^{-1}=\mathrm{Ob}\,X$, et tout revient à définir les
\[
\Hom_{X\Sigma^{-1}}(x,y),\qquad x,y\in\mathrm{Ob}\,X,
\]
\msp{page 6}
et pour $x,y,z\in\mathrm{Ob}\,X$ la composition
\[
\Hom_\Sigma(x,y)\times\Hom_\Sigma(y,z)\xrightarrow{(u,v)\mapsto vu}\Hom_\Sigma(x,z),
\]
(où on note $\Hom_\Sigma$ au lieu de $\Hom_{X\Sigma^{-1}}$, pour alléger les notations).

Grâce à $\Sigma$, on définit une sous-catégorie pleine
\[
\uCh_\Sigma(X)\hookrightarrow\uCh(X),
\]
(\footnote[3]{L'ensemble de ses objets est noté $\mathrm{Ch}(X)$. Ses éléments sont les $\Sigma$-chemins de $x$ à $y$ \texttt{[de} $x$ \texttt{à} $y$\texttt{?]}.}) par la condition sur $c\in\mathrm{Ch}(X)$ : Pour $i$ impair, \texttt{[plutôt pair]} i.e. $i-1\xleftarrow{u_i}i$ ($u_i$ est dite `rétrograde'), $c(u_i)\in\Sigma$, i.e. les `tronçons' rétrogrades dans le chemin $c$ doivent être dans $\Sigma$. On définit de même les sous-catégories pleines
\[
\uCh_\Sigma(X;x,y)=\{c\in\uCh(X;x,y)\mid c\in\uCh_\Sigma(X)\},
\]
ou encore la fibre en $(x,y)$ du foncteur source-but pour $\uCh_\Sigma(X)$
\[
\uCh_\Sigma(X)\longrightarrow X\times X.
\]

On définit une application canonique
\[
\mathrm{Ch}_\Sigma(X,x,y)\longrightarrow\Hom_\Sigma(x,y),
\]
\msp{page 7}
en associant `a tout chemin de $x$ dans $y$ le $\Sigma$-homomorphisme de $x$ dans $y$ obtenu en \textit{inversant} (dans $X\Sigma^{-1}$) les tronçons rétrogrades (ce qui est possible, puisqu'ils sont dans $\Sigma$), de sorte qu'on trouve dans $X\Sigma^{-1}$ un `chemin direct' % [sic] `a -> à, accent imprimé comme guillemet ouvrant
\[
x=A_{-n}\longrightarrow\cdots\longrightarrow A_{-1}\longrightarrow A_0\longrightarrow A_1\longrightarrow A_2\cdots A_n=A_{n+1}=\cdots=y.
\]
et on compose tous ces morphismes.

En fait, on trouve aussitôt que l'on définit ainsi un \textit{foncteur}
\[
\uCh_\Sigma(X)\longrightarrow\uFl(X\Sigma^{-1}),
\]
\texttt{[}$\uFl$ \texttt{catégorie des flèches]} du fait que si on a un diagramme

\printedpage{6}
% head: A. Grothendieck : Dérivateurs, Ch. V, Édition des Universités de Montpellier II et Paris VII
% foot: Version 9 septembre 2001
\[
\begin{tikzcd}
A_{i-1} \arrow[d, "f_{i-1}"] & A_i \arrow[l, "u_i"'] \arrow[d, "f_i"] \\
A'_{i-1} & A'_i \arrow[l, "u'_i"']
\end{tikzcd}
\qquad\qquad (i\text{ pair})
\]
avec $u_i$, $u'_i$ dans $\Sigma$, ce diagramme, en tant que diagramme dans $X\Sigma^{-1}$, est commutatif avec flèches horizontales des \textit{isomorphismes,} et reste donc commutatif en inversant lesdites flèches. D'où un foncteur induit
\[
\uCh_\Sigma(X;x,y)\longrightarrow\uFl(\underbrace{X\Sigma^{-1}}_{X'};x,y).
\]

Mais pour toute catégorie $X'$, et des objets $x$, $y$ dans $X'$, la catégorie $\uFl(X',x,y)$ est
\msp{page 8}
une catégorie \textit{discrète,} ayant comme ensemble d'objets $\Hom_{X'}(x,y)$. Ainsi le foncteur précédant définit en fait une application canonique % [sic] précédant -> précédent, adjectif
\[
\boxed{\pi_0\uCh_\Sigma(X;x,y)\longrightarrow\Hom_\Sigma(x,y).}
\]
Je dis que cette application est toujours bijective. Pour le voir, nous allons \textit{constuire} `from scratch' la catégorie $X\Sigma^{-1}$, en prenant comme objets ceux de $X$, et en \textit{définissant} % [sic] constuire -> construire, lettre omise
\[
\boxed{\Hom_\Sigma(x,y)=\pi_0\uCh_\Sigma(X;x,y).}
\]
Il faut donc définir la composition
\[
\Hom_\Sigma(x,y)\times\Hom_\Sigma(y,z)\longrightarrow\Hom_\Sigma(x,z).
\]
Pour ceci, on va définir des foncteurs `composition de chemins'
\[
\begin{aligned}
\uCh_m(X;x,y)\times\uCh_n(X;y,z) &\longrightarrow \uCh_{m+n}(X;x,z)\\
(c,c') &\longmapsto c'\circ c,
\end{aligned}
\]
mais en nous bornant pour simplifier au cas où $m$, $n$ ont \textit{la même parité,} de telle façon qu'en `mettant bout à bout' $c$ et $c'$ (dans cet ordre), on trouve bien un chemin de type $n+m$.

\msp{page 9}
Cette opération est `transposée' du couple de foncteurs
\[
\begin{aligned}
\Phi_n &\xrightarrow{\alpha} \Phi_{n+m}\quad i\longmapsto (i-m)\\
\Phi_m &\xrightarrow{\beta} \Phi_{n+m}\quad i\longmapsto (i+n),
\end{aligned}
\]
\printedpage{7}
% head: A. Grothendieck : Dérivateurs, Ch. V, Édition des Universités de Montpellier II et Paris VII
% foot: Version 9 septembre 2001
où on aura
\[
\alpha(\{n\})=\beta(\{-m\})=n-m,
\]
et on aura
\[
\Phi_{n+m}\simeq\Phi_n\sqcup_e\Phi_m,
\]
i.e. $\Phi_{n+m}$ est la `somme amalgamée' dans (Cat), obtenue en `recollant' $\Phi_n$ et $\Phi_m$, par identification de l'`élément but' de $\Phi_n$ avec l'`élément source' de $\Phi_m$. Le foncteur composition de chemins induit \texttt{[dans ce qui précède, à partir du début de la page 9, il faut inverser} $m$ \texttt{et} $n$\texttt{]}
\[
\uCh_{m,\Sigma}(x,y)\times\uCh_{n,\Sigma}(y,z)\longrightarrow\uCh_{m+n,\Sigma}(x,z),
\]
d'où en passant aux $\pi_0$
\[
\Hom_{\Sigma,m}(x,y)\times\Hom_{\Sigma,n}(y,z)\longrightarrow\Hom_{\Sigma,m+n}(x,z),
\]
où on dénote, pour $a,b\in\mathrm{Ob}\,X$, par
\[
\Hom_{\Sigma,n}(a,b)\overset{\text{déf}}{=}\pi_0(\uCh_{m,\Sigma}(a,b)).
\]

\texttt{[plutôt} $\uCh_{n,\Sigma}$\texttt{]}

\msp{page 10}
Pour en déduire une application sur les $\varinjlim$
\[
\underbrace{\Hom_\Sigma(x,y)}_{=\ \varinjlim\ \Hom_{\Sigma,m}(x,y)}\quad\times\quad\underbrace{\Hom_\Sigma(y,z)}_{=\ \varinjlim\ \Hom_{\Sigma,n}(y,z)}\quad\longrightarrow\quad\underbrace{\Hom_\Sigma(x,z)}_{=\ \varinjlim\ \Hom_{\Sigma,n+m}(x,z)}
\]
il faut prouver, pour $n\leq n'$, $m\leq m'$, la commutativité dans
\[
\begin{tikzcd}
\Hom_{\Sigma,m}(x,y)\times\Hom_{\Sigma,n}(y,z) \arrow[r] \arrow[d] & \Hom_{\Sigma,m+n}(x,y) \arrow[d] \\
\Hom_{\Sigma,m'}(x,y)\times\Hom_{\Sigma,n'}(y,z) \arrow[r] & \Hom_{\Sigma,m'+n'}(x,y)
\end{tikzcd}
\]
\texttt{[plutôt} $\Hom_{\Sigma,m+n}(x,z)$\texttt{,} $\Hom_{\Sigma,m'+n'}(x,z)$\texttt{]} En d'autres termes, que pour $c:x\longrightarrow y$ et $c':y\longrightarrow z$ deux $\Sigma$-chemins composables de longueur $m$ et $n$ respectivement, le composé $c'c:x\longrightarrow z$ ne change pas, `à homotopie près' dans $\uCh_\Sigma(x,z)$, quand on remplace $c$, $c'$ par des chemins $\bar{c}$, $\bar{c}'$ de longueur $m'$, $n'$ déduits de $c$ et $c'$ par `dégénérescence aux bouts'.

Plus précisement, % [sic] précisement -> précisément, accent omis
\[
\bar{c}'\bar{c}\sim\overline{c'c},
\]
le chemin composé est homotope à celui déduit de $c'c$ par `dégénérescence aux bouts'.

Je viens de regarder pour $n=m=1$, $n'=m'=2$, et de constater qu'il n'y a pas de flèche naturelle entre $\bar{c}'\bar{c}$ et $\overline{c'c}$ - l'homotopie cherchée doit
\printedpage{8}
% head: A. Grothendieck : Dérivateurs, Ch. V, Édition des Universités de Montpellier II et Paris VII
% foot: Version 9 septembre 2001
\msp{page 11}
se définir à l'aide d'un `chemin' dans $\uCh_{\Sigma,m',n'}(X;x,z)$. Dans le cas général, on va en fait trouver une homotopie
\[
h=h_{(n,m),(n',m')}
\]
`rendant commutatif' le diagramme de foncteurs
\[
\begin{tikzcd}[column sep=large]
\mathcal{A}=\uCh_{\Sigma,m}(X;x,y)\times\uCh_{\Sigma,n}(X;y,z) \arrow[r, "*_{m,n}"] \arrow[d, "i_{m,m'}\times j_{n,n'}"] & \uCh_{\Sigma,m+n}(X;x,z) \arrow[d, "k_{m+n,m'+n'}", ""'{name=K}] \\
\uCh_{\Sigma,m'}(X;x,y)\times\uCh_{\Sigma,n'}(X;y,z) \arrow[r, "*_{m',n'}", ""{name=S}] & \uCh_{\Sigma,m'+n'}(X;x,z)=\mathcal{B},
\arrow[from=K, to=S, bend right, "h"']
\end{tikzcd}
\]
où les $i_{\alpha,\beta}$, $j_{\alpha,\beta}$, $k_{\alpha,\beta}$ sont les foncteurs de transition
\[
\uCh_{\Sigma,\alpha}(X;a,b)\longrightarrow\uCh_{\Sigma,\beta}(X;a,b)
\]
$(\alpha\leq\beta)$ pour les $(x,y)$, $(y,z)$ et $(x,z)$ (en guise de $(a,b)$) respectivement. L'homotopie $h$ sera un chemin dans la catégorie $\uHom(\mathcal{A},\mathcal{B})$, allant de
\[
k_{m+n,m'+n'}\circ{*}_{m+n}\longrightarrow{*}_{m'+n'}\circ(i_{m,m'}\times j_{n,n'}).
\]
L'existence d'une telle homotopie signifie ni plus ni moins que le diagramme carré ci-dessus est commutatif, en tant que diagramme dans la catégorie (Cat\,hot) des `catégories à homotopie près,' dont les objets sont
\msp{page 12}
les (petites) catégories, et où on définit
\[
\Hom_{\mathrm{Cat\,hot}}(\mathcal{A},\mathcal{B})\overset{\text{déf}}{=}\underbrace{\pi_0(\uHom(\mathcal{A},\mathcal{B}))}_{\begin{array}{@{}l@{}}\text{classes d'homotopie}\\ \text{de foncteurs de }\mathcal{A}\\ \text{dans }\mathcal{B}\end{array}}\,.
\]
C'est cette conclusion qui sera pour nous utile, et non la description explicite du chemin $h$. Cette existence admise (cf. ci-dessous la preuve), on voit donc que si on regarde, pour $a,b\in\mathrm{Ob}\,X$, le système des catégories $\uCh_{\Sigma,r}(X;a,b)$ $(r\in\mathbf{N})$ comme un \textit{ind-objet de} Cat\,hot\textit{,} qu'on désigne par
\[
\uCh^{\mathrm{ind}}_{\Sigma}\,(X;a,b)\in\mathrm{Ob}\,\mathrm{Ind}(\mathrm{Cat\,hot}),
\]
on aura un morphisme canonique dans Ind(Cat\,hot)
\[
\uCh^{\mathrm{ind}}_{\Sigma}\,(X;x,y)\times\uCh^{\mathrm{ind}}_{\Sigma}\,(X;y,z)\xrightarrow{\ *\ }\uCh^{\mathrm{ind}}_{\Sigma}\,(X;x,z),
\]
appelé \textit{`ind-morphisme de composition de chemins.'} On aurait aimé le déduire d'un honnête \textit{foncteur} `composition de chemins'
\[
\begin{aligned}
\uCh_\Sigma(X;x,y)\times\uCh_\Sigma(X;y,z) &\xrightarrow{\ ?\ } \uCh_\Sigma(X;x,z),\\
\text{qu'on noterait }(c,c') &\longmapsto c'\circ c\text{ ou }c'\cdot c,
\end{aligned}
\]
mais je doute qu'un tel foncteur existe,

% ---- from Der-ChV-234_p011-015.tex
\printedpage{9}
% head: A. Grothendieck : Dérivateurs, Ch. V, Édition des Universités de Montpellier II et Paris VII

\msp{page 13}
qui par passage aux ind-objets redonne le morphisme précédent, ou, ce qui revient au même, tel que pour tout $m,n\in\mathbf{N}$ (de même parité), il existe un $N\geq m+n$ assez grand tel que l'on ait un diagramme
\[
\begin{tikzcd}
\Ch_{\Sigma,m}(X;x,y)\times\Ch_{\Sigma,n}(X;y,z) \arrow[r,"*_{m,n}"] \arrow[dd,hook] \arrow[dr] & \Ch_{\Sigma,m+n}(X;x,z) \arrow[d,hook] \\
 & \Ch_{\Sigma,N}(X;x,z) \arrow[d,hook] \\
\Ch_{\Sigma}(X;x,y)\times\Ch_{\Sigma}(X;y,z) \arrow[r,"?"] & \Ch_{\Sigma}(X;x,z)
\end{tikzcd}
\]
et où le triangle supérieur est \textit{commutatif dans} Cat\,hot. (Il faut supposer que pour $N$ assez grand, le composé
\begin{tikzcd}[cramped, sep=small]
{} \arrow[d,hook] & \\
{} \arrow[r,"?"] & {}
\end{tikzcd}
se factorise par $\Ch_{\Sigma,N}(X;x,z)$ \dots)
Pour voir que les foncteurs composition $*_{m,n}$ définissent un homomorphisme des ind-objets dans (Cat\,hot), il suffit de le prouver dans les cas particuliers $(m',n')=(m+1,n+1)$. On va faire la figure en supposant $m$, $n$ \textit{impairs.}

Je m'aperçois ici que je me suis canulé stupidement, la composition de deux chemins de longueur \textit{impaire $m$, $n$ n'est}

\msp{page 14}
\textit{pas un chemin} de longueur $m+n$, i.e. un diagramme de type $\Phi_n$, mais un diagramme de type \textit{dual} $\Phi_n^o=\Psi_n$. Par exemple, pour $m=n=1$
\[
\underbrace{A_0\longleftarrow A_1\longrightarrow A_2}_{c}=\underbrace{B_0\longleftarrow B_1\longrightarrow B_2}_{c'}
\]
est de type $\Phi_2^o$! Donc pour ne pas trop compliquer la vie, on va dans les développements précédents se borner à $m$, $n$ \textit{pairs,} donc des chemins de longueur au moins 4. C'est un peu ridicule! On est ramené à prouver la commutativité à homotopie près dans le cas où on laisse $c'$ fixé et où on envoie $\Ch_{\Sigma,m}(X;x,y)$ dans $\Ch_{\Sigma,m+2}(X;x,y)$, ou inversement $c$ fixe et $\Ch_{\Sigma,n}(X;y,z)$ dans $\Ch_{\Sigma,n+2}$ \texttt{[lire:}$\Ch_{\Sigma,n+2}(X;y,z)$\texttt{]}. Prenons le premier cas, et notons un chemin $c$ sans indices négatifs (ce n'est pas une notation bien heureuse!)
\[
A_0\longrightarrow A_1\longleftarrow A_2\cdots\longleftarrow A_{2m}
\]
(NB à cause de la condition $n$ pair, tous nos chemins commencent désormais par une flèche de $A_0$ dans $A_1$, jamais en sens inverse).

Je vois que l'explicitation de l'homotopie, sans offrir de difficulté, est très fastidieuse à expliquer. Je vais essayer de donner une démonstration conceptuelle, sans `calcul' fastidieux.

% foot: Version 9 septembre 2001

\printedpage{10}
% head: A. Grothendieck : Dérivateurs, Ch. V, Édition des Universités de Montpellier II et Paris VII

\msp{page 15}
Les deux foncteurs composés
\[
f=k\circ *_{m,n}\ ,\quad g=*_{m',n'}\circ(i\times j)
\]
dans le carré page 11, dont il faut prouver qu'ils sont homotopes dans $\Hom(\mathcal{A},\mathcal{B})$, peuvent être vus comme obtenus de la façon suivante, en termes de deux foncteurs
\[
\varphi,\psi:\Phi_{m'+n'}\rightrightarrows\Phi_{m+n}, \tag*{(1)}
\]
d'où deux foncteurs `transposés'
\[
\begin{tikzcd}
\Hom(\Phi_{m,n},X) \arrow[r,shift left,"F"] \arrow[r,shift right,"G"'] & \Hom(\Phi_{m',n'},X)
\end{tikzcd} \tag*{(2)}
\]
\texttt{[plutôt} $\Phi_{m+n}$\texttt{,} $\Phi_{m'+n'}$\texttt{]} lesquels foncteurs `induisent' $f$ et $g$. De façon précise, désignant par $a$, $b$ les objets source et but de $\Phi_{m+n}$, par $c$ l'objet qui correspond au but du chemin $c$ de longueur $m$ \texttt{[trop de} $c$\texttt{'s]}, à la source du chemin $c'$ de longueur $n$ qu'il s'agit de composer, enfin par $a_0$, $b_0$ les éléments source et but de $\Phi_{m'+n'}$, les foncteurs envisagés (1) appliquent $a_0$ dans $a$, $b_0$ dans $b$, donc définissent un diagramme commutatif
\[
\begin{tikzcd}
\Phi_{m'+n'} \arrow[r,shift left,"\varphi"] \arrow[r,shift right,"\psi"'] & \Phi_{m+n} \\
\underset{\text{\footnotesize catégorie discrète}}{\varepsilon=\{a_0,b_0\}.} \arrow[u,hook,"\alpha'"] \arrow[ur,hook,"\alpha"] &
\end{tikzcd} \tag*{(3)}
\]
Ceci dit, la catégorie but $\mathcal{B}$ pour $f$, $g$, à savoir $\Ch_{\Sigma,m'+n'}(X;x,z)$, s'interprète comme la sous-catégorie pleine de $\Hom(\Phi_{m+n},X)$ formée

\msp{page 16}
des foncteurs $c':\Phi_{m'+n'}\longrightarrow X$ dont le composé avec $\alpha'$ est le foncteur
\[
\beta:\varepsilon\longrightarrow X\qquad\qquad
\begin{aligned}
a_0&\longrightarrow x\\
b_0&\longrightarrow z,
\end{aligned} \tag*{$(*)$}
\]
et de même la catégorie source $\mathcal{A}$ s'interprète comme la sous-catégorie pleine de $\Hom(\Phi_{m+n},X)$ formée des $c:\Phi_{m+n}\longrightarrow X$ dont la restriction à $\{a,b,c\}$ est le foncteur
\[
a\longrightarrow x,\ b\longrightarrow y,\ c\longrightarrow z.
\]
\texttt{[plutôt} $b\longrightarrow z$\texttt{,} $c\longrightarrow y$\texttt{]} Elle est contenue dans la sous-catégorie pleine $\mathcal{A}'$ des foncteurs dont la restriction à $\{a,b\}$ soit $a\longrightarrow x,\ b\longrightarrow z$, i.e. dont le composé avec $\alpha$ soit le foncteur $\beta:\varepsilon\longrightarrow X$ déjà envisagé ci-dessus. Ainsi les foncteurs $f,g:\mathcal{A}\rightrightarrows\mathcal{B}$ se factorisent par $\mathcal{A}'$
% foot: Version 9 septembre 2001
\printedpage{11}
% head: A. Grothendieck : Dérivateurs, Ch. V, Édition des Universités de Montpellier II et Paris VII
\[
\begin{tikzcd}
\mathcal{A} \arrow[r,hook] \arrow[dr,shift left,"f"] \arrow[dr,shift right,"g"'] & \mathcal{A}' \arrow[d,shift right,"g'"'] \arrow[d,shift left,"f'"] \\
 & \mathcal{B},
\end{tikzcd}
\]
foncteurs déduits de $\varphi$, $\psi$ ci-dessus, donnant lieu aux triangles commutatifs (3), et du foncteur $\beta:\varepsilon\longrightarrow X$, de la façon dite. Pour prouver que $f$ et $g$ sont homotopes, il suffit de le prouver pour $f'$, $g'$. Et pour ceci, il suffit de montrer que $\varphi$ et $\psi$ sont homotopes en tant qu'objets de la catégorie $\Hom(\Phi_{m+n},\Phi_{m'+n'};\alpha)$ formée des foncteurs dont la restriction à $\varepsilon$ est $\alpha$ \texttt{[plutôt de la catégorie} $\Hom(\Phi_{m'+n'},\Phi_{m+n};\alpha')$ \texttt{formée des foncteurs dont la restriction à} $\varepsilon$ \texttt{est} $\alpha'$ \texttt{]}.

\msp{page 17}
D'ailleurs, les deux assertions (sur $f'$, $g'$, quels que soient $X$, $x$, $y$, $z$, \texttt{[plutôt} $x$\texttt{,} $z$\texttt{]} et sur $\varphi$, $\psi$, $\alpha$ \texttt{[plutôt} $\alpha'$\texttt{]}) sont en fait équivalentes comme on voit en faisant $X=\Phi_{m+n}$, $x=a$, $y=c$, $z=b$ \texttt{[plutôt} $x=a$\texttt{,} $z=b$\texttt{]}. Ceci sera contenu dans le lemme plus général. (\footnote[4]{Dans les considérations qui suivent, on considère en fait tacitement que les foncteurs envisagés sont \textit{croissants} sur les sommets pour l'ordre de leur numérotation. Donc il faudrait écrire $\Hom_{cr}$ au lieu de $\Hom$, dans les énoncés qui suivent, y compris dans le lemme de connexité ci-contre.})

\textbf{Lemme.} \textit{Soient $k$ et $l$ deux entiers. Alors la sous-catégorie pleine de $\Hom(\Phi_k,\Phi_l)$ formée des foncteurs $\Phi_k\longrightarrow\Phi_l$ appliquant source en source, but en but, est connexe.} (NB Elle est non vide si et seulement si $k\geq l$.)

On va décrire entièrement ces foncteurs (= applications croissantes) $\Phi_k\xrightarrow{\varphi}\Phi_l$, en identifiant $\Fl(\Phi_k)$ \texttt{[plutôt l'ensemble des flèches non identiques de} $\Phi_k$\texttt{]} à l'intervalle d'entiers $[1,2k]$, et en associant à $\varphi$ la partie
\[
E_\varphi\subset\Fl(\Phi_k)=[1,2k]\qquad\qquad E_\varphi=\{u\in\Fl(\Phi_k)\mid\varphi(u)\text{ une identité}\}
\]
formée des $u\in\Fl(\Phi_k)$ qui sont `contractées' par $\varphi$ en une identité. Si $E$ est un ensemble d'entiers, on y introduit la notion de composantes connexes (par la relation d'équivalence dans $E$ engendrée par la relation de succession $y=x+1$). On voit aussitôt que les composantes connnexes de $E_\varphi$ qui ne contiennent ni la première ni la dernière flèche de $\Phi_k$ sont paires, et que % [sic] connnexes -> connexes, triple n
\[
\underbrace{\card\Fl(\Phi_k)}_{2k}=\card\,(E_\varphi)+\underbrace{\card\Fl(\Phi_l)}_{2l},
\]

\msp{page 18}
i.e.
\[
\card\,(E_\varphi)=2(k-l).
\]

% foot: Version 9 septembre 2001

\printedpage{12}
% head: A. Grothendieck : Dérivateurs, Ch. V, Édition des Universités de Montpellier II et Paris VII

Il est immédiat qu'un foncteur $\varphi$ est surjectif sur les objets et sur les flèches, donc identifie $\Phi_l$ à une catégorie quotient de $\Phi_k$ - et la relation d'équivalence dans $\Phi_k$ définie par $\varphi$ est celle obtenue en identifiant tous les objets qui sont source ou but de flèches d'une même composante connexe $I$ de $E_\varphi$
\[
\cdots\longleftarrow\ \underbrace{\longrightarrow\ \longleftarrow\ \longrightarrow\ \longleftarrow}_{I}\ \longrightarrow\cdots
\]
et les flèches dans cette composante à l'identité correspondante. Les conditions énoncées sur $E$ (indépendamment d'un $\varphi$ qui lui donnerait naissance assurent que la catégorie quotient est isomorphe (avec \textit{préservation} de source et de but) soit à $\Phi_l$, soit à la duale $\Psi_l=\Phi_l^o$, suivant la direction de la première flèche de $\bar{\Phi}_k$ (quotient de $\Phi_k$). Cette `direction initiale' est la même que celle pour $\Phi_k$ si la première flèche de $\Phi_k$ n'appartient pas à $E$, ou bien si elle y appartient et sa composante connexe dans $E$ est de cardinal pair. Dans le cas inverse ($u\in E$ et sa composante connnexe dans $E$ est de cardinal impair), la direction initiale de $\bar{\Phi}_k$ est opposée à celle de $\Phi_k$. % [sic] "(indépendamment d'un $\varphi$ qui lui donnerait naissance" -> parenthèse fermante manquante, the opening parenthesis is never closed; [sic] connnexe -> connexe, triple n

\msp{page 19}
On voit donc que $E$ doit satisfaire encore à une condition supplémentaire, pour que le $\bar{\Phi}_k$ correspondant soit isomorphe à $\Phi_l$, i.e. pour que $E$ provienne d'un foncteur $\varphi$. Je résume toutes les conditions

\begin{itemize}
\item[a)] $\card E=2(k-l)$.
\item[b)] Les composantes connexes de $E$ qui ne contiennent pas de flèche initiale ou finale sont de cardinal pair.
\item[c)] Si $\Phi_k$ et $\Phi_l$ ont même direction initiale ($k$ et $l$ de même parité), alors l'éventuelle composante de $E$ contenant la flèche initiale doit être elle aussi de cardinal pair (ce qui, à cause de a) \texttt{[plutôt a) et b)]}, équivaut à la même condition pour l'éventuelle composante connexe contenant la flèche finale de $\Phi_k$). Si $\Phi_k$ et $\Phi_l$ ont des directions initiales différentes, il faut que la flèche initiale de $\Phi_k$ appartienne à une composante connexe de cardinal \textit{impair} de $E$ (ce qui équivaut à itou pour la flèche finale).
\end{itemize}

D'autre part, le foncteur $\varphi$ est uniquement déterminé par la connaissance de $E$, car il y a un \textit{unique} isomorphisme $\bar{\Phi}_k\xrightarrow{\sim}\Phi_l$ compatible avec sources et buts, à cause du

\textbf{Lemme} (immédiat) : \textit{Un automorphisme de $\Phi_l$ qui est l'identité sur l'objet source (ou sur l'objet but) est l'identité partout.}

Les foncteurs $\varphi$, $\psi$ de tantôt correspondaient au choix, l'un de $E_\varphi=\{1,2,k-1,k\}$ \texttt{[plutôt} $E_\varphi=\{1,2,2k-1,2k\}$\texttt{]}, l'autre de $E_\psi=\{1,2,i,i+1\}$, 2 composantes connexes

\msp{page 20}
dans un cas, trois \texttt{[plutôt deux]} dans l'autre.) % [sic] "l'autre.)" -> parenthèse fermante sans parenthèse ouvrante, no matching opening parenthesis

% foot: Version 9 septembre 2001

\printedpage{13}
% head: A. Grothendieck : Dérivateurs, Ch. V, Édition des Universités de Montpellier II et Paris VII

Dans le cas général, il y a lieu d'expliciter la relation d'ordre de l'ensemble $\Hom^{!}(\Phi_k,\Phi_l)$ (! signifie : compatible avec source et but), en termes des ensembles représentatifs $E_\varphi$, $E_\psi$. Je ne vais pas expliciter cette relation d'ordre, mais seulement une propriété de celle-ci qui nous sera suffisante. Soit $\varphi_E$ le foncteur défini par une partie $E\subset[1,2k]$. Tout d'abord :

\textbf{Lemme 1} (immédiat) : \textit{Soient $E,E'\subset[1,2k]$ admissibles, et soit $j\in[1,2k]$ tel que $E\cap[1,j]=E'\cap[1,j]$. Alors $\varphi_E$ et $\varphi_{E'}$ coïncident sur le but de $u_j$ ($j$ une flèche de $\Phi_k$) et sur tous les objets de $\Phi_k$ antérieurs à celle-ci (au sens de la numérotation des objets, et non de la relation d'ordre de la catégorie $\Phi_k$ bien sûr.) Dualement, si $E\cap[j,2k]=E'\cap[j,2k]$, alors $\varphi_E$ et $\varphi_{E'}$ coïncident sur la source de $y$, et sur tous les objets de $\Phi_k$ postérieurs à celui-ci.} % [sic] la source de $y$ -> la source de $u_j$, apparently intended by duality with "le but de $u_j$"

Considérons maintenant une partie connexe $I$ de $E$ (pas forcément une composante connexe), $I=[r,s]\subset[1,2k]$ ($r$ le plus petit élément de $I$, $s$ son plus grand élément) et supposons que

\begin{itemize}
\item[a)] $\card I$ est pair (i.e. $s-r$ est pair) \texttt{[plutôt} $s-r$ \texttt{impair]}
\item[b)] $r\neq1$, et $r-1\notin E$ (i.e. $r-1\in[1,2k]\backslash E$)
\end{itemize}

et soit
\[
E'=(E\backslash\{s\})\cup\{r-1\}=E\backslash I\cup\overbrace{[r-1,s-1]}^{I'}
\]
déduit de $I$ en `translatant' sur la gauche d'un cran.

\msp{page 21}
Il est immédiat que si $E$ est admissible, $E'$ aussi (et inversement). Ceci dit, on a le

\textbf{Lemme 2.} \textit{$E$ et $E'$ étant comme ci-dessus, il existe une homomorphisme de foncteurs entre $\varphi_E$ et $\varphi_{E'}$} (je ne dis pas dans quel sens va la flèche!). % [sic] une homomorphisme -> un homomorphisme, wrong gender

\textbf{Corollaire.} Énoncé dual, quand $E'$ est déduit de $E$ en translatant à droite d'un cran un intervalle connexe de longueur paire $I=[r,s]$ de $E$, $s\neq2k$ et $s+1\notin E$.

En effet, dans ce cas $E$ est déduit de $E'$ par une opération du type explicité dans le lemme 2.

La connexité de la catégorie $\Hom^{!}(\Phi_k,\Phi_l)$ est maintenant claire. Si $k$ et $l$ ont même parité, en poussant les composantes connexes de $E$ vers la gauche, on écrit $E$ sous la forme `normale'
\[
E_0=[1,2(k-l)].
\]
Si $k$ et $l$ sont de parité inverse, on ne peut pousser toute la composante connexe de $E$ qui contient l'élément dernier $2k$, car elle est de cardinal impair. Il faut laisser $2k$ dans $E$ et pousser le reste, et on ramène $E$ à la forme normale
\[
E_0=[1,2(k-l)-1]\cup\{2k\}.
\]

\hfill cqfd!

% foot: Version 9 septembre 2001

% ---- from Der-ChV-234_p016-020.tex
\printedpage{14}
% head: A. Grothendieck : Dérivateurs, Ch. V, Édition des Universités de Montpellier II et Paris VII
Ce qui précède justifie notre raisonnement, qui donnait par composition des chemins des applications de composition bien définies
\[
\underbrace{\Hom_\Sigma(x,y)}_{=\pi_0\uCh_\Sigma(x,y)} \times \underbrace{\Hom_\Sigma(y,z)}_{=\pi_0\uCh_\Sigma(y,z)} \longrightarrow \underbrace{\Hom_\Sigma(x,z)}_{=\pi_0\underline{\textstyle\mathrm{Ch}}_\Sigma(x,z)}.
\]
% (the third "Ch" under the brace is printed in a larger size than the two others)

\msp{page 22}Je dis que cette composition est associative. Ceci résulte du fait que si on a trois chemins (d'indices \textit{pairs}!)
\[
x \xrightarrow{c} y \xrightarrow{c'} z \xrightarrow{c''} t,
\]
on a (sans avoir à effectuer des homotopies \dots) associativité stricte
\[
c''(c'c) = (c''c')c.
\]

Je dis qu'il y a des identités. On définit
\[
1_x \in \Hom_\Sigma(x,x)
\]
par le chemin d'indice 0 (qui est pair!), i.e. par le diagramme réduit à $x$, sans flèche (ou si on veut interpréter le diagramme comme un foncteur $e \longrightarrow X$, il y a la flèche identique sous-entendue). Nos raisonnements n'excluaient pas le cas $n = 0$ ou $m = 0$, et il est tautologique que pour un chemin
\[
x \xrightarrow{c} y
\]
on a
\[
c \cdot 1_x = 1_y \cdot c = c.
\]
Donc on a bien défini une catégorie, que je vais appeler $X_\Sigma$. Ja vais définir un foncteur canonique % [sic] Ja -> Je, misprint of the pronoun
\[
\theta : X \longrightarrow X_\Sigma
\]
qui sera l'identité sur les objets, il faut le définir sur les flèches
\[
\Hom_X(x,y) \longrightarrow \Hom_{X_\Sigma}(x,y) = \pi_0\mathrm{Ch}_\Sigma(x,y)
\]
\texttt{[plutôt} $\pi_0\uCh_\Sigma(x,y)$\texttt{]}. Malheureusement, une flèche
\[
x \xrightarrow{u} y
\]
\msp{page 23}ne peut s'interpréter directement comme un chemin à notre sens, il faut payer nos fuites! Mais on va à une telle flèche associer le chemin
\[
x \longrightarrow y \xleftarrow{\mathrm{id}} y \xrightarrow{\mathrm{id}} y \xleftarrow{\mathrm{id}} y,
\]
% foot: Version 9 septembre 2001
\printedpage{15}
% head: A. Grothendieck : Dérivateurs, Ch. V, Édition des Universités de Montpellier II et Paris VII
élément de $\uCh_{\Sigma,2}(x,y)$, ce qui définit l'application cherchée. Il faut vérifier la compatibilité avec les identités et sous la composition. À la flèche $1_x : x \longrightarrow x$ est associé l'élément de $\Phi_2(x,x)$
\[
x \xrightarrow{\mathrm{id}} x \xleftarrow{\mathrm{id}} x \xrightarrow{\mathrm{id}} x \xleftarrow{\mathrm{id}} x,
\]
il faut prouver que dans $\pi_0(\uCh_{\Sigma,2}(x,x))$ il devient égal à l'élément déduit de $1_x \in \Phi_{\Sigma,0}(x,x)$, par le foncteur de transition. Mais cet élément n'est autre justement que
\[
x \xrightarrow{\mathrm{id}} x \xleftarrow{\mathrm{id}} x \xrightarrow{\mathrm{id}} x \xleftarrow{\mathrm{id}} x,
\]
c'est le même (sans avoir même à passer au $\pi_0$)!

Pour la compatibilité avec la composition
\[
x \xrightarrow{u} y \xrightarrow{v} z,
\]
on a d'une part
\[
x \xrightarrow{vu} z \xleftarrow{\mathrm{id}} z \xrightarrow{\mathrm{id}} z \xleftarrow{\mathrm{id}} z,
\]
d'autre part
\[
x \xrightarrow{u} y \xleftarrow{\mathrm{id}} y \xrightarrow{\mathrm{id}} y \xleftarrow{\mathrm{id}} y \xrightarrow{v} z \xleftarrow{\mathrm{id}} z \xrightarrow{\mathrm{id}} z \xleftarrow{\mathrm{id}} z.
\]
On complète le premier chemin par dégénérescence aux deux bouts, pour le comparer au deuxième en tant qu'élément de $\uCh_4(x,z)$,
\[
x \xrightarrow{\mathrm{id}} x \xleftarrow{\mathrm{id}} x \xrightarrow{vu} z \xleftarrow{\mathrm{id}} z \xrightarrow{\mathrm{id}} z \xleftarrow{\mathrm{id}} z \xrightarrow{\mathrm{id}} z \xleftarrow{\mathrm{id}} z.
\]

\msp{page 24}Il faut prouver que les deux sont homotopes. On a donc deux foncteurs
\[
f,g : \underbrace{\overbrace{\uFl(X;x,y)}^{\text{catégorie discrète}} \times \overbrace{\uFl(X;y,z)}^{\text{catégorie discrète}}}_{\mathcal{A}} \underset{g}{\overset{f}{\rightrightarrows}} \underbrace{\uCh_{4,\Sigma}(X;x,z)}_{\mathcal{B}}, \tag*{$(\ast)$}
\]
on veut prouver qu'ils sont homotopes. Pour ceci, on les regarde (comme tantôt) comme `transposés' des deux foncteurs
\[
\begin{tikzcd}[column sep=large]
\bDelta_1 \sqcup \bDelta_1 \simeq \bDelta_2 & \Phi_4 \arrow[l, shift right, "\varphi"'] \arrow[l, shift left, "\psi"] \\
& \varepsilon, \arrow[ul, "\alpha"] \arrow[u, "\alpha'"']
\end{tikzcd}
\]
vu que les catégories $\mathcal{A}$, $\mathcal{B}$ des deux membres de (*) s'interprètent comme des sous-catégories pleines de $\uHom(\bDelta_2, X)$ et $\uHom(\Phi_4, X)$ respectivement; à savoir les sous-catégories pleines formées des foncteurs dont le composé avec $\alpha$ resp. $\alpha'$ est donné par $x$ et par $z$, et de plus (pour $\mathcal{A}$) dont la valeur en l'objet médian de $\bDelta_2$ soit $y$. On peut encore remplacer $\mathcal{A}$ par la sous-catégorie pleine plus grande $\mathcal{A}'$ de $\uHom(\bDelta_2, X)$, obtenue en ignorant cette dernière condition, et il suffit de prouver que les deux foncteurs `plus gros'
% foot: Version 9 septembre 2001
\printedpage{16}
% head: A. Grothendieck : Dérivateurs, Ch. V, Édition des Universités de Montpellier II et Paris VII
\[
\mathcal{A}' \underset{g'}{\overset{f'}{\rightrightarrows}} \mathcal{B}'
\]
sont homotopes. Et pour ceci, il suffit de
\msp{page 25}voir encore que $\varphi$ et $\psi$ sont homotopes dans la catégories des foncteurs $\Phi_4 \longrightarrow \bDelta_2$ qui sont compatibles avec la source, et le but. On voit donc que notre lemme d'homotopie de tantôt n'était pas assez général. Il faudrait définir les \textit{catégories intervalles généralisés,} de longueur $n$ quelconque ($n =$ nombre de flèches bout à bout), une catégorie intervalle généralisé \textit{standard} de longueur $n$ a pour ensemble d'objets l'intervalle d'entiers % [sic] la catégories -> la catégorie, singular article with plural noun
\[
[0,n],
\]
et pour tout couple d'entiers consécutifs $i-1$, $i$ parmi les objets ($1 \leq i \leq n$), il y a exactement une flèche qui relie $i$, $i+1$ \texttt{[plutôt} $i-1$\texttt{,} $i$\texttt{]}. On affecte l'intervalle $[i,i+1]$ du signe $+$ si $i \longrightarrow i+1$ (flèche en sens direct), du signe $i$ \texttt{[plutôt} $-$\texttt{]} si $i \longleftarrow i+1$ (flèche en sens rétrograde) - donc les catégories standards sont définies, pour $n$ fixé, par un système de signes $+$ ou $-$, dont on affecte les $n$ `tronçons' consécutifs $A_i := [i-1,i]$ ($1 \leq i \leq n$) de l'intervalle d'entiers $[0,n]$. Il y en a donc exactement $2^n$.

On s'attend à ce que l'énoncé suivant soit vrai :

\textbf{Lemme.} \textit{Soient $I$, $I'$ deux intervalles généralisés, munis d'objets source et objet but (de sorte que $I$, $I'$ soient} canoniquement \textit{isomorphes aux intervalles genéralisés standards correspondants). Soit} $\uHom^!(I,I')$ % [sic] genéralisés -> généralisés, missing accent on the first e
\msp{page 26}\textit{la catégorie des foncteurs de $I$ dans $I'$, compatibles avec source et but. Si cette catégorie est $\neq \emptyset$, elle est} 0\textit{-connexe} (et, je présume, même contractile?).

Je vais admettre provisoirement qu'il en est ainsi, du moins dans le cas particulier de $\Phi_4$ (intervalle généralisé de longueur 8) et $\bDelta_2$. Cela donne donc un foncteur
\[
\theta : X \longrightarrow X_\Sigma.
\]
On a gagné, car la propriété universelle de ce foncteur est pratiquement évidente. Mais j'ai oublié de vérifier que $\theta$ transforme éléments de $\Sigma$ en flèches inversibles. Soit donc
\[
x \xrightarrow{u} y
\]
avec $u \in \Sigma$, on va définir un chemin en sens inverse
\[
y \xrightarrow{c} x
\]
qui définira dans $X_\Sigma$ une flèche inverse de $u$. L'idée naturelle est de prendre le `chemin'
\[
y \xleftarrow{u} x
\]
% foot: Version 9 septembre 2001
\printedpage{17}
% head: A. Grothendieck : Dérivateurs, Ch. V, Édition des Universités de Montpellier II et Paris VII
de $y$ dans $x$, mais ce n'est pas un chemin au sens de nos définitions, on va le compléter en
\[
c : y \xrightarrow{\mathrm{id}} y \xleftarrow{u} x \xrightarrow{\mathrm{id}} x \xleftarrow{\mathrm{id}} x.
\]
Je dis que $c$ fait l'affaire. On doit prouver deux relations d'homotopie, de $c \cdot \bar{u}$ et de $\bar{u} \cdot c$ avec les chemins identiques en $x$ \texttt{[resp. en} $y$\texttt{]} de longueur 8 (indice 4, i.e. dans $\uCh_{\Sigma,4}(X;x,x)$ resp. $\uCh_{\Sigma,4}(X;y,y)$),
\msp{page 27}où $\bar{u}$ est le chemin associé à $u$
\[
x \xrightarrow{u} y \xleftarrow{\mathrm{id}} y \xrightarrow{\mathrm{id}} y \xleftarrow{\mathrm{id}} y.
\]

On a donc à regarder
\[
x \xrightarrow{u} y \xleftarrow{\mathrm{id}} y \xrightarrow{\mathrm{id}} y \xleftarrow{\mathrm{id}} y \xrightarrow{\mathrm{id}} y \xleftarrow{u} x \xrightarrow{\mathrm{id}} x \xleftarrow{\mathrm{id}} x
\]
et voir qu'il est homotope au chemin trivial
\[
x \xrightarrow{\mathrm{id}} x \xleftarrow{\mathrm{id}} \cdots x \xrightarrow{\mathrm{id}} x \xleftarrow{\mathrm{id}} x
\]
(des identités partout), et itou pour
\[
y \xrightarrow{\mathrm{id}} y \xleftarrow{u} x \xrightarrow{\mathrm{id}} x \xleftarrow{\mathrm{id}} x \xrightarrow{u} y \xleftarrow{\mathrm{id}} y \xrightarrow{\mathrm{id}} y \xleftarrow{\mathrm{id}} y
\]
et le chemin trivial en $y$. Cette fois, on a encore deux foncteurs
\[
\begin{tikzcd}[column sep=large]
\bDelta_1 & \Phi_4 \arrow[l, shift right, "\varphi"'] \arrow[l, shift left, "\psi"] \\
& \varepsilon, \arrow[ul, "\alpha'"] \arrow[u, "\alpha"']
\end{tikzcd}
\]
compatibles avec objets sources et objets buts, et on doit prouver qu'ils sont homotopes dans $\uHom^!(\Phi_4, \bDelta_1)$. Cela doit donc résulter du même lemme général. On a gagné - par une interprétation simple de $X\Sigma^{-1}$, et en particulier des $\Hom_{X\Sigma^{-1}}(x,y)$ et de leur composition, en termes de la géométrie des espaces de chemins :
\[
\left\{
\begin{aligned}
&\Hom_{X\Sigma^{-1}} \simeq \pi_0(\uCh_\Sigma(X;x,y)) \\ % [sic] Hom_{XΣ^{-1}} -> Hom_{XΣ^{-1}}(x,y), arguments (x,y) missing
&\text{Composition des homomorphismes dans } X\Sigma^{-1} \text{ correspond à} \\
&\text{la composition des } \Sigma\text{-chemins dans } X
\end{aligned}
\right.
\]
% foot: Version 9 septembre 2001

\printedpage{18}
% head: A. Grothendieck : Dérivateurs, Ch. V, Édition des Universités de Montpellier II et Paris VII
\noindent\msp{page 28}

\section*{3\quad Passage aux catégories de diagrammes}

Nous aurions dû noter la catégorie $X$ $\mathcal{M}$, car par la suite elle jouera le rôle d'une `catégorie de modèles'. On désignera par contre par $X$ les `catégories d'indices' sur lesquelles on veut construire les `catégories de diagrammes' de type $X$
\[
\mathcal{M}(X) = \uHom(X,\mathcal{M}).
\]
On se donne donc
\[
\underbrace{\Sigma}_{\text{quasi-isomorphismes}} \subseteq \mathrm{Fl}(\mathcal{M}), \qquad \Sigma \supset \text{ensemble des flèches identiques},
\]
et pour toute catégorie $X$ on définit
\[
\Sigma_X \subset \mathrm{Fl}\,\mathcal{M}(X)
\]
ainsi: si $F, G \in \mathrm{Ob}\,\mathcal{M}(X)$ ($F, G : X \rightrightarrows \mathcal{M}$), une flèche $u : F \longrightarrow G$ (homomorphisme de foncteurs) est dans $\Sigma_X$ si et seulement si $\forall x \in X$, $u(x) : F(x) \longrightarrow G(x)$ est dans $\Sigma$ (`critère de quasi-isomorphie fibre par fibre', pour l'homomophisme entre les copréfaisceaux $F$, $G$ envisagés). % [sic] homomophisme -> homomorphisme, letter r missing

Soient
\[
a, b \in \mathrm{Ob}\,\mathcal{M}
\]
et considérons les foncteurs constants correspondants sur $X$
\[
a_X, b_X \in \mathrm{Ob}\,\mathcal{M}(X).
\]

\msp{page 29}On se propose d'appliquer la théorie du \S2 au `calcul' de
\[
\Hom_{\Sigma_X}(a_X, b_X),
\]
on a donc une bijection canonique
\[
\boxed{\Hom_{\Sigma_X}(a_X, b_X) \simeq \pi_0\uCh_{\Sigma_X}(\mathcal{M}(X); a_x, b_x).} % [sic] a_x, b_x -> a_X, b_X, lowercase index against a_X, b_X elsewhere
\]
On se propose d'expliciter la catégorie de chemins du second membre. On a, abstraction faite des données de $\Sigma$, $a$, $b$ dans $\mathcal{M}$, l'expression pour la catégorie de \textit{tous} les chemins
% foot: Version 9 septembre 2001

% ---- from Der-ChV-234_p021-025.tex
\printedpage{19}
% head: A. Grothendieck : Dérivateurs, Ch. V, Édition des Universités de Montpellier II et Paris VII

dans $\cM(X)$
\[
\begin{aligned}
\Ch(\cM(X)) \quad &\simeq \quad \Hom(\ \underbrace{\boldsymbol{\Phi}}_{\substack{\text{pro-catégorie}\\ (\Phi_n)_{n\in\mathbf{N}}}}\ ,\ \underbrace{\cM(X)}_{\Hom(X,M)}\ ) \\ % ?? uncertain: M in "Hom(X,M)" under the brace -- italic M or calligraphic \mathcal{M}
&\dfeq \quad \varinjlim_n \Hom(\Phi_n,\cM(X)) \\
&\simeq \quad \varinjlim_n \Hom(\Phi_n\times X,\cM) \\
&\simeq \quad \varinjlim_n \Hom(X,\Hom(\Phi_n,\cM)) \\
&\dfeq \quad \varinjlim_n \Hom(X,\Ch_n(\cM)) \\
&\simeq \quad \Hom^{!}(X,\Ch(\cM)),
\end{aligned}
\]
où le signe ! signifie qu'on prend la sous-catégorie pleine des foncteurs
\[
X\longrightarrow \Ch(\cM),
\]
\msp{page 30}
qui se factorisent par une des sous-catégories pleines $\Ch_n(\cM)$. On peut dire aussi que si on considère le système inductif des catégories $\Ch_n(\cM)$ comme un ind-objet de Cat, qu'on va noter
\[
\Chh(\cM) \qquad\qquad\qquad\qquad \text{ind-objet de Cat,}
\]
on a
\[
\Ch(\cM(X)) \xrightarrow{\ \sim\ } \Hom(X,\Chh(\cM)).
\]
Si on en vient à $\Sigma$, $\Sigma_X$, on trouve aussitôt qu'un chemin dans $\cM(X)$ est un $\Sigma_X$-chemin, si et seulement si le foncteur correspondant $X\longrightarrow\Ch(\cM)$ se factorise par $\Ch_{\Sigma}(\cM)$. Ainsi on trouve
\[
\Ch_{\Sigma_X}(\cM(X))\simeq\Hom^{!}(X,\Ch_{\Sigma}(\cM)),
\]
où le signe ! a le même sens que plus haut. Si on introduit le ind-objet $\Chh_{\Sigma}(\cM)$ défini par les $\Ch_{\Sigma,n}(\cM)$, on trouve donc aussi
\[
\Ch_{\Sigma_X}(\cM(X))\xrightarrow{\ \sim\ }\Hom(X,\Chh_{\Sigma}(\cM)).
\]

\msp{page 31}
Soit $c$ un chemin dans $\cM(X)$, et
\[
X\longrightarrow\Ch(\cM)
\]
le foncteur qu'il définit, composons-le avec les foncteurs source-but
\[
\Ch(\cM)\longrightarrow\cM\times\cM,
\]
on trouve un foncteur
\[
X\longrightarrow\cM\times\cM,
\]
i.e. un objet $(F_0,F_1)$ de $\cM(X)\times\cM(X)$. Je dis que $F_0$ et $F_1$ sont respectivement la source et le but du chemin $c$. De façon un peu plus précise, on a commutativité (stricte) dans le diagramme de foncteurs
% foot: Version 9 septembre 2001
\printedpage{20}
% head: A. Grothendieck : Dérivateurs, Ch. V, Édition des Universités de Montpellier II et Paris VII
\[
\begin{tikzcd}[column sep=large, row sep=large]
\Ch_{\Sigma_X}(\cM(X)) \arrow[r, "\sim"] \arrow[d, "{\sigma_{\cM(X)},\,\beta_{\cM(X)}}"] & \Hom^{!}(X,\Ch_{\Sigma}(\cM)) \arrow[d, "{\text{déduit de }\sigma_{\cM},\,\beta_{\cM}}"] \\
\cM(X)\times\cM(X) \arrow[r, "\sim"] & \Hom(X,\cM\times\cM),
\end{tikzcd}
\]
où les foncteurs $\sigma$ sont les foncteurs source, les foncteurs $\beta$ les foncteurs but.

Si on se donne arbitrairement les objets $F_0$, $F_1$ dans $\cM(X)$, donc un foncteur
\[
X\longrightarrow\cM\times\cM,
\]
$\Ch_{\Sigma_X}(F_0,F_1)$ s'identifie donc à la catégorie des foncteurs ‘admissibles’ (i.e. se factorisant par un $\Ch_{\Sigma,n}(\cM)$) $X\xrightarrow{\ C\ }\Ch_{\Sigma}(\cM)$ qui \textit{relèvent}
\msp{page 32 vide}
\msp{page 33}
le foncteur $(\sigma_{\cM},\beta_{\cM})$ \texttt{[plutôt} $(F_0,F_1)$\texttt{]} précédent, i.e. qui donnent lieu à un triangle commutatif (strict)
\[
\begin{tikzcd}[row sep=large]
 & \Ch_{\Sigma}(\cM) \arrow[d, "{(\sigma_{\cM},\,\beta_{\cM})}"] \\
X \arrow[ur, "C"] \arrow[r, "{(F_0,F_1)}"] & \cM\times\cM.
\end{tikzcd}
\]
Mais si les foncteurs $F_0$, $F_1$ sont constants, de valeur $a$, $b$ comme tantôt, i.e. $(F_0,F_1)$ constant de valeur $(a,b)$, alors la catégorie des foncteurs $C$ en question est tautologiquement isomorphe à celle des foncteurs de $X$ dans la fibre de $(\sigma_{\cM},\beta_{\cM})$ en $(a,b)$, d'où l'isomorphisme de catégories fondamental
\[
\boxed{\Ch_{\Sigma_X}(X;a_X,b_X)\simeq\Hom^{!}(X,\Ch_{\Sigma}(\cM;a,b)).}
\]

On va expliciter le contenu essentiel de cette formule en une

\textbf{Scholie.} Soient $\cM$ une catégorie avec un ensemble $\Sigma$ de flèches de localisation, $a,b\in\mathrm{Ob}\,\cM$. Alors on peut trouver un objet $\KK=K(a,b)$, $K=(K_n)$, de Ind\,Cat (indexé par $\mathbf{N}$ mais peu importe ce fait), et un isomorphisme canonique pour toute catégorie $X$ (isomorphisme fonctoriel en $X$)
\[
\begin{aligned}
\mathrm{Hom}_{\Sigma_X}(a_X,b_X) \quad &\simeq \quad \pi_0(\Hom(X,\KK)) \\
&\dfeq \quad \pi_0(\ \varinjlim \Hom(X,K_n)) \\
&\simeq \quad \varinjlim \pi_0(\Hom(X,K_n)).
\end{aligned}
\]
% foot: Version 9 septembre 2001

\printedpage{21}
% head: A. Grothendieck : Dérivateurs, Ch. V, Édition des Universités de Montpellier II et Paris VII

\msp{page 34}
Je vais épiloguer sur ce résultat, sans me soucier pour l'instant de la façon particulière dont la ind-catégorie $K$ \texttt{[plutôt} $\KK$\texttt{]} a été obtenue en termes de $\cM$, $\Sigma$, $a$, $b$.

Comme les $\pi_0\Hom(X,K_n)=\mathrm{Hom}_{\Cathot}(X,K_n)$ ne dépendent que de $X$ et de $K_n$ en tant qu'objets de Cat\,hot, je vais interpréter $\KK$ comme un ind-objet de Cat\,hot, et écrire la formule fondamentale sous la forme
\[
\mathrm{Hom}_{\Sigma_X}(a_X,b_X)\simeq\mathrm{Hom}_{\mathrm{Ind\,Cat\,hot}}(X,\KK).
\]

Ainsi, l'énoncé du scholie implique deux choses :

\begin{enumerate}
\item[a)] Le contrafoncteur $X\longmapsto\mathrm{Hom}_{\Sigma_X}(a_X,b_X)$ sur Cat se factorise canoniquement par Cat\,hot, via un foncteur canonique
\[
\Cathot^{o}\ \xrightarrow{\ H=H^{\cM,\Sigma}_{a,b}\ }\mathrm{Ens}\ ;
\]
\item[b)] Ce dernier contrafoncteur sur Cat\,hot est \textit{ind-représentable} par un ind-objet $\KK$ de Cat\,hot.
\end{enumerate}

C'est cet ind-objet de Cat\,hot qui est défini à isomorphisme unique près par $\cM$, $\Sigma$, $a$, $b$, indépendamment de tout artifice de calcul pour expliciter les $\mathrm{Hom}_{\Sigma_X}(a_X,b_X)$.

Notons aussi la conclusion suivante, nullement évidente à priori :

\msp{page 35}
Notons d'abord (ceci étant implicite dans le scholie) que toute flèche dans Cat
\[
X\xrightarrow{\ f\ }Y
\]
définit
\[
\cM(Y)\xrightarrow{\cM(f)}\cM(X) \qquad\qquad\qquad \text{appliquant } \Sigma_Y \text{ dans } \Sigma_X,
\]
d'où
\[
\cM(Y)\Sigma_Y^{-1}\dfeq\mathbf{D}_{\cM,\Sigma}(Y)\longrightarrow\cM_X\Sigma_X^{-1}=\mathbf{D}_{\cM,\Sigma}(X),
\]
\texttt{[plutôt} $\cM(X)$\texttt{]} d'où, pour tout $a,b\in\mathrm{Ob}\,\cM$, vu que $\cM(f)$ transforme $a_Y$, $b_Y$ en $a_X$, $b_X$ :
\[
\boxed{\mathrm{Hom}_{\Sigma_Y}(a_Y,b_Y)\longrightarrow\mathrm{Hom}_{\Sigma_X}(a_X,b_X).}
\]

Ceci posé, on a :

\textbf{Corollaire :} \textit{Si $f$ est un homotopisme, alors quelque soient $a,b\in\cM$, la flèche précédente est un isomorphisme.} % [sic] quelque soient -> quels que soient, grammatical form of the locution

La question naturelle, c'est de savoir si la même conclusion est encore valable, en supposant seulement que $f$ soit un ‘\textit{quasi-isomorphisme}’ (ou ‘équivalence faible’), ici que $f$ définisse un isomorphisme dans la catégorie de fractions Hot de Cat. (À ne pas confondre avec Cat\,hot,
% foot: Version 9 septembre 2001
\printedpage{22}
% head: A. Grothendieck : Dérivateurs, Ch. V, Édition des Universités de Montpellier II et Paris VII
\msp{page 36}
nettement plus fine). S'il en était ainsi, le contrafoncteur
\[
X\longmapsto\mathrm{Hom}_{\Sigma_X}(a_X,b_X)
\]
se factoriserait non seulement par Cat\,hot, mais même par la catégorie de fractions plus grossière Hot. Et dès lors se poserait la question si ce contrafoncteur hypothétique
\[
\mathrm{Hot}^{o}\longrightarrow\mathrm{Ens}
\]
est encore ind-représentable; et de préférence (\footnote[5]{ce sera automatique si on a ind-représentabilité comme foncteur sur Hot, cf. ci-dessous.}) ind-représentable par le ind-objet image de $\KK$ dans Ind\,Hot par
\[
\mathrm{Ind}(\Cathot)\longrightarrow\mathrm{Ind}(\mathrm{Hot}),
\]
déduit du foncteur de localisation $\Cathot\longrightarrow\mathrm{Hot}$. (Car cet ind-objet est un candidat tout trouvé.)

\msp{page 37}
Je vais faire quelques rappels catégoriques, généralisant la situation du foncteur canonique de localisation
\[
\Cathot\xrightarrow{\ \varphi\ }\mathrm{Hot},
\]
dont on sait qu'il admet un adjoint à droite \texttt{[si un lecteur a une démonstration de ce fait, prière de contacter G. Maltsiniotis : [e-mail address removed]]}
\[
\Cathot\xleftarrow{\ \psi\ }\mathrm{Hot}, \qquad\qquad\qquad \mathrm{Hom}_{\mathrm{Hot}}(\varphi X,Y)\simeq\mathrm{Hom}_{\Cathot}(X,\psi Y),
\]
où $\psi$ est le foncteur ‘enveloppe de Kan’ ou ‘représentant de Kan’. On va donc considérer la situation plus générale
\[
\begin{tikzcd}[column sep=large]
C \arrow[r, shift right=1, "\varphi"'] & C' \arrow[l, shift right=1, "\psi"']
\end{tikzcd}
\qquad\qquad \mathrm{Hom}_{C'}(\varphi X,X')\simeq\mathrm{Hom}_{C}(X,\psi X')
\]
i.e.
\[
\varphi,\psi \text{ couple de foncteurs adjoints.}
\]
On suppose de plus
\[
\left\{
\begin{aligned}
&\varphi \text{ un foncteur localisation, i.e. } \psi \text{ pleinement fidèle,}\\
&\text{i.e. } \varphi\psi\xrightarrow[\sim]{\text{adj}}\mathrm{id}_{C'} \quad .
\end{aligned}
\right.
\]
Soit $W=W_{\varphi}\subset\mathrm{Fl}\,C$, défini par $\varphi$ : $W_{\varphi}=\{u\in\mathrm{Fl}\,C\mid\varphi(u)\text{ isomorphisme}\}$. Voici la proposition que j'avais oubliée, et que j'ai eu du mal à reconstituer, hélas!

\textbf{Proposition 1.} \textit{Soit $K\in \mathit{Ob}C$. Les conditions suivantes sur $K$ sont équivalentes :}
% foot: Version 9 septembre 2001

\printedpage{23}
% head: A. Grothendieck : Dérivateurs, Ch. V, Édition des Universités de Montpellier II et Paris VII

\begin{enumerate}
\item[(i)] \textit{Le foncteur}
\[
X\longmapsto K(X)\dfeq\mathrm{Hom}_C(X,K) \qquad\quad C^{o}\longrightarrow\mathrm{Ens}
\]
\textit{transforme flèches de $W$ (‘quasi-isomorphismes’) en bijections.}
\item[(i$'$)] \textit{Même condition pour les flèches d'adjonction}
\[
X\longrightarrow\psi\varphi X \qquad\quad (\forall X\in\mathrm{Ob}\ \mathrm{C})
\]
\item[(i$''$)] \textit{Même condition pour la flèche d'adjonction} (\footnote[6]{N.B. Il suffit même que l'élément $\mathrm{id}_K$ de $F_K(K)$ provienne de $F_K(\psi\varphi K)$ \texttt{[où} $F_K=K(-)=\mathrm{Hom}(-,K)$\texttt{]}, i.e.\par (i$'''$) Il existe une rétraction de $\psi\varphi X$ sur $X$ \texttt{[plutôt} $\psi\varphi K$ \texttt{sur} $K$\texttt{]} (NB Celle-ci sera même un isomorphisme).})
\[
K\longrightarrow\psi\varphi K.
\]
\item[(ii)] \textit{$K$ appartient à l'image essentielle de $\psi$}
\item[(ii$'$)] \textit{$K\longrightarrow\psi\varphi K$ (adjonction) est un isomorphisme.}
\item[(ii$''$)] \textit{L'homomorphisme fonctoriel en $X$}
\[
\mathrm{Hom}_C(X,K)\longrightarrow\mathrm{Hom}_{C'}(\varphi X,\varphi K)
\]
\textit{est un isomorphisme.}
\end{enumerate}

\textsc{Démonstration.} On a d'abord comme quasi-tautologies (ii$'$)$\Longrightarrow$ (ii)$\Longrightarrow$ (ii$''$)$\Longrightarrow$ (ii$'$), i.e. l'équivalence des trois dernières conditions. En effet, (ii’)$\Longrightarrow$ (ii) tautologique, (ii)$\Longrightarrow$ (ii$''$) car si $K=\psi(K')$ on a % ?? uncertain: the primes in "(ii’)" after "En effet," and in "(ii'')" -- text apostrophe or math prime
\msp{page 38}
par adjonction
\[
\mathrm{Hom}_C(X,K)=\mathrm{Hom}_C(X,\psi(K'))\overset{\text{adj}}{\simeq}\mathrm{Hom}_{C'}(\varphi X,K'),
\]
or $\varphi(K)=\varphi\psi(K')\xrightarrow{\ \sim\ }K'$ (adjonction, et $\psi$ pleinement fidèle).
\[
\begin{tikzcd}[column sep=small, row sep=large]
\mathrm{Hom}_C(X,K) \arrow[rr, "\sim"] \arrow[dr, "\varphi_{X,K}"] & & \mathrm{Hom}_C(\varphi X,K') & \text{\footnotesize qui se complète en} \\ % [sic] Hom_C(\varphi X,K') -> Hom_{C'}(\varphi X,K'), \varphi X and K' are objects of C' (cf. the display above)
& \mathrm{Hom}_{C'}(\varphi X,\ \underbrace{\varphi K}_{=\varphi\psi(K')}\ ) \arrow[ur, "\sim", "\text{adj}"'] & &
\end{tikzcd}
\]
ce qui prouve que $\varphi_{X,K}$ est un isomorphisme.

(ii$''$) $\Longrightarrow$ (ii$'$). Car la flèche
% foot: Version 9 septembre 2001

% ---- from Der-ChV-234_p026-030.tex
\printedpage{24}
% head: A. Grothendieck : Dérivateurs, Ch. V, Édition des Universités de Montpellier II et Paris VII
% foot: Version 9 septembre 2001

\[
\begin{tikzcd}
\Hom_C(X,K) \arrow[r,"\alpha_X"] \arrow[dr,dashed] & \Hom_{C'}(\varphi X,\varphi K) \arrow[d,dash,"\wr"',"\mathrm{adj}"] \\
 & \Hom_C(X,\psi\varphi K)
\end{tikzcd}
\]
est déduite de la flèche d'adjonction
\[
K \longrightarrow \psi\varphi K
\]
en prenant $\Hom_C(X,-)$, donc $\alpha_X$ est un isomorphisme pour tout $X$ si et seulement si cette flèche d'adjonction est un isomorphisme. Donc
\[
\text{(ii)} \Longleftrightarrow \text{(ii$'$)} \Longleftrightarrow \text{(ii$''$)}.
\]
D'autre part,
\[
\text{(ii$''$)} \Longrightarrow \text{(i)} \Longrightarrow \text{(i$'$)} \Longrightarrow \text{(i$''$)}
\]
sont tautologiques, en tenant compte pour (i) $\Longrightarrow$ (i) \texttt{[probablement (i) $\Longrightarrow$ (i$'$)]} que les flèches d'adjonction
\[
X \xrightarrow{u} \psi\varphi X \qquad\qquad\qquad \text{dans } C
\]
sont bien dans $W$, car on a un triangle commutatif
\[
\begin{tikzcd}[column sep=large]
\varphi(X) \arrow[r,"\varphi(u)"] \arrow[dr,"\id_{\varphi(X)}"] & \varphi(\psi\varphi X) = (\varphi\psi)(\varphi X) \arrow[d,"{\begin{array}{l}\wr\ \mathrm{adj}\ \varphi\psi \longrightarrow 1,\ \text{qui est un}\\ \text{isomorphisme car } \psi \text{ est pleinement fidèle}\end{array}}"] \\
 & \varphi(X),
\end{tikzcd}
\]
donc $\varphi(u)$ est bien un isomorphisme.

Reste à prouver (i$''$) $\Longrightarrow$ (ii$'$), i.e. on suppose

\noindent\msp{page 39}%
que le foncteur $X \longrightarrow \Hom(X,K)$ transforme le morphisme d'adjonction particulier
\[
K \xrightarrow{i} \psi\varphi K
\]
en
\[
\begin{aligned}
u \quad &\longrightarrow \quad u\circ i\\
\Hom(\psi\varphi(K),K) \quad &\longrightarrow \quad \Hom(K,K),
\end{aligned}
\]
qui est une bijection. Donc $\id_K$ est dans l'image, soit $p$ un élément du premier membre qui lui donne naissance

\printedpage{25}
% head: A. Grothendieck : Dérivateurs, Ch. V, Édition des Universités de Montpellier II et Paris VII
% foot: Version 9 septembre 2001

\[
\begin{tikzcd}
K \arrow[r,hook,shift left,"i"] & \psi\,\underbrace{\varphi(K)}_{K'} \arrow[l,shift left,"p"]
\end{tikzcd}
\qquad\qquad {\scriptstyle \text{donc } p\circ i = \id_K}
\]
Considérons
\[
ip = q : \psi(K') \longrightarrow \psi(K') \qquad\qquad\qquad (\text{NB } q^2 = q),
\]
comme $\psi$ est pleinement fidèle, il existe un unique
\[
q_0 : K' \longrightarrow K' \qquad\qquad\qquad (K' \overset{\text{déf}}{=} \varphi(K))
\]
tel que l'on ait
\[
q = \psi(q_0).
\]
La relation $pi = \id_K$ donne, en appliquant $\varphi$,
\[
\varphi(p)\varphi(i) = \id,
\]
donc comme $\varphi(i)$ est un isomorphisme (c'est
\[
\begin{tikzcd}[column sep=large]
\varphi(K) \arrow[r,"\sim"] \arrow[dr,"\sim"{sloped,pos=0.45},"\id"{pos=0.6}] & \varphi(\psi\varphi)(K) = (\varphi\psi)(\varphi K) \arrow[d,"\wr"',"\mathrm{adj}"] \\
 & \varphi(K),
\end{tikzcd}
\]
rendant commutatif le triangle ci-contre), $\varphi(p)$ l'est. Donc $\varphi(q) = \varphi(i)\varphi(p)$ l'est, puisque $\varphi(i)$ et $\varphi(i)$ \texttt{[plutôt $\varphi(p)$]} le sont. Mais
\[
\varphi(q) = \varphi\psi(q_0) \quad \underbrace{\xrightarrow{\ \sim\ }}_{\mathrm{adj}\ \varphi\psi\ \xrightarrow{\sim}\ \id} \quad q_0,
\]
donc $q_0$ est un isomorphisme, donc aussi $q = \psi q_0$. Comme $q^2 = q$, cela implique $q = \id$. Donc $i$ et $p$ sont des isomorphismes inverses l'un de l'autre, ce qui établit (ii$''$) \texttt{[plutôt (ii$'$)]} et achève la démonstration.

\textbf{Corollaire :} \textit{Considérons des foncteurs $C^o \xrightarrow{H} \Ens$, et $C'^o \xrightarrow{H'} \Ens$, tels que $H \simeq H'\circ\varphi$. Pour que $H$ soit représentable, il faut et il suffit que $H'$ le soit, et dans ce cas, l'objet $K$ de $C$}

\noindent\msp{page 40}%
\textit{qui représente $H$ satisfait aux conditions équivalentes de la proposition. Si $K$, $K'$ les représentent, on a des isomorphismes canoniques $K' \simeq \varphi K$, $K \simeq \psi K'$.}

Si $H'$ est représentable par $K' \in \Ob C'$, on aura donc
\[
H'\circ\varphi(X) \simeq \Hom_{C'}(\varphi X,K') \overset{\mathrm{adj}}{\simeq} \Hom_C(X,\psi(K'))
\]
\printedpage{26}
% head: A. Grothendieck : Dérivateurs, Ch. V, Édition des Universités de Montpellier II et Paris VII
% foot: Version 9 septembre 2001
donc $H$ est représentable par $K = \psi(K')$, satisfaisant la condition (ii) de la proposition. Inversement, si $H$ est représentable par un objet $K$, alors $K$ satisfait à la condition (i), donc par (ii$''$) on a
\[
\Hom_C(X,K) \xrightarrow{\sim} \Hom_{C'}(\varphi X,\varphi K),
\]
donc le foncteur $H' : C^o \longrightarrow \Ens$ qui factorise $H$ s'explicite comme $X' \longmapsto \Hom_{C'}(X',\varphi K)$, donc est représentable par $\varphi K$. % [sic] C^o -> C'^o, H' is a functor on C'^o (cf. the Corollaire)

Passons aux ind-objets. Les foncteurs $\varphi$, $\psi$ passent aux ind-objets, d'où des foncteurs
\[
\begin{tikzcd}[column sep=large]
\Ind C \arrow[r,shift left,"\Psi = \Ind\psi"] \arrow[r,shift right,"\Phi = \Ind\varphi"'] & \Ind C'.
\end{tikzcd} % [sic] both arrows printed pointing to Ind C' -> the arrow Psi = Ind psi apparently intended from Ind C' to Ind C, psi goes from C' to C
\]
On aura encore la formule d'adjonction pour $X = (X_i)_{i\in I}$ dans $\Ind C$, $X' = (X'_{i'})_{i'\in I'}$ dans $\Ind C'$, car
\[
\begin{aligned}
\Hom(X,\Psi X') &= \varprojlim_i \varinjlim_{i'} \Hom(X_i,\psi(X'_{i'}))\\
&= \varprojlim_i \varinjlim_{i'} \Hom(\varphi X_i,X'_{i'})\\
&= \Hom(\Phi X,X').
\end{aligned}
\]

\noindent\msp{page 41}%
On prouve de même que $\Psi$ est pleinement fidèle, car si $X' = (X'_i)_{i\in I}$, $Y' = (Y'_j)_{j\in J}$ sont des objets de $\Ind C'$, on aura
\[
\begin{tikzcd}[column sep=small]
\Hom(X',Y') \arrow[d,"\Psi_{X',Y'}"] & \simeq \varprojlim\limits_i \varinjlim\limits_j & \Hom_{C'}(X'_i,Y'_j) \arrow[d,"\psi_{X'_i,Y'_j}"] \\
\Hom(\Psi X',\Psi Y') & \simeq \varprojlim\limits_i \varinjlim\limits_j & \Hom_C(\psi(X'_i),\psi(Y'_j)),
\end{tikzcd}
\]
donc $\Psi_{X',Y'}$ est un isomorphisme si les $\psi_{X'_i,Y'_j}$ le sont.

Donc la proposition 1 peut s'appliquer à la situation $\Phi$, $\Psi$ entre ind-objets. Je vais l'expliciter sous la forme la plus commode :

\textbf{Proposition 2.} \textit{Soit $K = (K_i)_{i\in I}$ un ind-objet de $C$, et soit $H$ le foncteur}
\[
\left\{
\begin{aligned}
H : X \quad &\longmapsto \quad \underbrace{\Hom_C(X,K)}_{K(X)} \overset{\text{déf}}{=} \varinjlim_i \Hom_C(X,K_i)\\
C^o \quad &\longrightarrow \quad \Ens
\end{aligned}
\right.
\]
\textit{qu'il ind-représente. Conditions équivalentes}

\begin{enumerate}[label={}]
\item[(i)] \textit{Le foncteur $H$ se factorise (à isomorphisme près) par $C'$ via $\varphi$, i.e. transforme les $u \in W_\varphi$ en isomorphismes.}
\printedpage{27}
% head: A. Grothendieck : Dérivateurs, Ch. V, Édition des Universités de Montpellier II et Paris VII
% foot: Version 9 septembre 2001
\item[(i$'$)] \textit{Même condition pour les flèches d'adjonction $X \longrightarrow \psi\varphi(X)$ (lesquelles sont dans $W$).}
\item[(i$''$)] \textit{Même condition pour la flèche d'adjonction $K \longrightarrow \psi\varphi(K)$.}
\item[(ii)] \textit{$K$ appartient à l'image essentielle de} $\Ind C'$ \textit{par $\Psi$.}
\item[(ii$'$)] $K \longrightarrow \Psi\Phi(K)$ \textit{est un} isomorphisme \textit{dans} $\Ind C$.) % [sic] ".)" -> ".", closing parenthesis without opening parenthesis
\item[(ii$''$)] \textit{L'homomorphisme fonctoriel}
\[
K(X) = \Hom_C(X,K) \longrightarrow \Hom_{C'}(\varphi(X),\varphi(K))
\]
\texttt{[plutôt} $\Hom_{\mathrm{Ind}(C)}(X,K) \longrightarrow \Hom_{\mathrm{Ind}(C')}(\Phi(X),\Phi(K))$\texttt{, en identifiant l'objet} $X$ \texttt{de} $C$ \texttt{à son image dans} $\mathrm{Ind}(C)$\texttt{]} \textit{est un isomorphisme.}
\end{enumerate}

\textsc{Démonstration.} L'équivalence de (ii), (ii$'$), (ii$''$) est la proposition 1 appliquée à $\Phi$ et $\Psi$, à cela près qu'il faut noter que (ii$''$) reste évidemment valable si on y remplace $X$ par un ind-objet $(X_j)_{j\in J}$ quelconque. On a

\noindent\msp{page 42}%
\[
\text{(ii$''$)} \Longrightarrow \text{(i)} \Longrightarrow \text{(i$'$)}
\]
tautologiquement, il reste à prouver que (i$'$) implique (i$''$) et (i$''$) (ii$'$). On ne peut pas appliquer directement la proposition 1, car l'hypothèse (i$'$) ici ne porte que sur les objets de $\Ind C$ qui proviennent de $C$, et il n'est pas évident à priori qu'elle passe aux objets quelconques. Mais cela l'est dans le cas d'un homomorphisme de ind-objets $(X_i)$, $(Y_i)$ sur le même ensemble d'indices, donné par un homomorphisme ordinaire (terme à terme) de systèmes inductifs sur $I$, $X_i \longrightarrow Y_i$: si les $X_i \longrightarrow Y_i$ sont dans $W$, l'hypothèse (i) implique que $H(Y) \longrightarrow H(X)$ est un isomorphisme, et de même pour (i$'$) pour $\underbrace{X}_{=(X_i)_i} \;\longrightarrow\; \underbrace{\Psi\Phi(X)}_{=(\psi\varphi(X_i))_i}$\,, où l'homomorphisme d'adjonction est donné terme à terme par les homomorphismes d'adjonction $X_i \longrightarrow \psi\varphi(X_i)$. Donc l'hypothèse (i$'$) ici implique l'hypothèse (i$'$) de la proposition 1 pour $\Phi$ et $\Psi$, donc aussi (ii$'$) (par la proposition 1), et on gagne. % [sic] "et (i'') (ii')" -> "et (i'') => (ii')", the implication sign or verb is missing

Revenons à la situation où $C = \Cathot$, $C' = \Hot$, $\varphi$ le foncteur \texttt{[de]} localisation. Si $K \in \Ob\Cathot$, alors $\underline{\psi\varphi(K)}$ s'appelle \textit{l'enveloppe de Thomason-Kan} de $K$ - si $K$ est une catégorie, son enveloppe $\bar{K}$ et la flèche $K \longrightarrow \bar{K}$ ne sont définis que dans la catégorie Cat\,hot, i.e. $\bar{K}$ n'est défini qu'à `homotopisme près'. % ?? uncertain: underline under psi varphi(K) -- underlined / not underlined

\noindent\msp{page 43}%
Mais on peut toujours \textit{choisir} un objet $\bar{K}$ de Cat, et une flèche
\[
i : K \longrightarrow \bar{K} \qquad\qquad\qquad \textit{dans } \Cat,
\]
\printedpage{28}
% head: A. Grothendieck : Dérivateurs, Ch. V, Édition des Universités de Montpellier II et Paris VII
% foot: Version 9 septembre 2001
pas seulement dans Cat\,hot, de telle façon de plus que $i$ soit une `\textit{cofibration} de Thomason' (et en particulier une immersion ouverte?). Il en résulte que si on a un système inductif
\[
K = (K_n)_{n\in\mathbf{N}}
\]
indexé par $\mathbf{N}$, et si on choisit pour chaque $n$ une enveloppe de Thomason-Kan
\[
K_n \overset{\alpha_n}{\hookrightarrow} \bar{K}_n
\]
de cette façon, on peut trouver des foncteurs de transition
\[
\bar{K}_n \xrightarrow{\bar{i}_n} \bar{K}_{n+1}
\]
tels que les carrés
\[
\begin{tikzcd}
K_n \arrow[r,hook,"\alpha_n"] \arrow[d,"i_n"'] & \bar{K}_n \arrow[d,"\bar{i}_n"] \\
K_{n+1} \arrow[r,hook,"\alpha_{n+1}"] & \bar{K}_{n+1}
\end{tikzcd}
\]
soient commutatifs, de telle sorte que l'on ait donc un homomorphisme de \textit{systèmes inductifs de} Cat
\[
K \longrightarrow \bar{K}
\]
et a fortiori un homomorphisme de ind-objets de Cat. Ceci posé, les conditions (i), (ii$'$) de la proposition 2 montrent que le foncteur ind-représenté par $K$ \textit{dans} $C = \Cathot$ transforme quasi-isomorphismes

\noindent\msp{page 44}%
ordinaires (équivalences faibles) en bijections, \textit{si et seulement si} l'homomorphisme précédent est un isomorphisme de ind-objets de Ind\,Cat\,hot.

Ce sera donc cela la condition naturelle à imposer aux ind-catégories de chemins
\[
\underline{\underline{\Ch}}_{\Sigma}(\mathcal{M};a,b) \qquad\qquad\qquad (a,b \in \mathrm{ob}\,\mathcal{M})
\]
provenant d'une `catégorie de modèles' $\mathcal{M}$, $\Sigma$.

Bien sûr, disposant d'un système inductif explicite de catégories (pas seulement d'objets de Cat\,hot)
\[
K_n = \underline{\Ch}_{\Sigma,n}(\mathcal{M};a,b),
\]
il est tentant de passer à la limite inductive
\[
K_\infty = \varinjlim_n \underline{\Ch}_{\Sigma,n}(\mathcal{M};a,b) \overset{\text{déf}}{=} \underline{\Ch}_{\Sigma}(\mathcal{M};a,b),
\]
et ceci d'autant plus que les foncteurs de transition sont des monomorphismes et pleinement fidèles, de sorte que les $K_n$ s'identifient à une suite de sous-catégories pleines de

% ---- from Der-ChV-234_p031-034.tex
\printedpage{29}
% head: A. Grothendieck : Dérivateurs, Ch. V, Édition des Universités de Montpellier II et Paris VII
% foot: Version 9 septembre 2001
$K_\infty$, dont la réunion est $K_\infty$. Mais d'autre part je ne sais donner aucune interprétation convaincante de $K_\infty$, comme de celle du ind-objet de Cat\,hot defini par $\uuK = (K_n)_{n\in\mathbf{N}}$, qui ind-représente le foncteur remarquable % [sic] defini -> défini, missing accent
\[
H = H_{\mathcal{M};x,y} : \Cathot^{o} \longrightarrow \Ens
\]
factorisant le foncteur
\msp{page 45}
\[
X \longmapsto H_{\mathcal{M};x,y}(X) \defeq \Hom_{\mathcal{M}(X)\Sigma_X^{-1}}(x_X, y_X),
\]
par la formule donc
\[
H(X) = \varinjlim_{n} \underbrace{\pi_0\uHom(X, K_n)}_{\Hom_{\Cathot}(X,K_n)} .
\]
Je ne sais donner d'interprétation de ce type ni de la catégorie $K_\infty$ elle-même, ni de son image dans Cat\,hot, ni de son image dans Hot. Il est vrai que si je restreins le foncteur $H$ à la sous-catégorie pleine de (Cat) formée des catégories \textit{finies} (Ob\,$X$ et Fl\,$X$ finis), alors on aura
\[
\uHom^{!}(X, K_\infty) = \uHom(X, K_\infty),
\]
et l'isomorphisme précédant s'interprète comme % [sic] précédant -> précédent, adjective intended
\[
\Hom_{\mathcal{M}(X)\Sigma_X^{-1}}(a_X, b_Y) \defeq H(X) \simeq \pi_0\uHom(X, K_\infty) \defeq \Hom_{\Cathot}(X, K_\infty).
\]
\texttt{[plutôt} $(a_X, b_X)$\texttt{]} Cela prouve du moins que la connaissance de $K_\infty$ est utile. Mais cette formule, pour $X$ restreint aux catégories \textit{finies,} ne définira pas $K_\infty$ comme objet de Cat\,hot, à isomorphisme unique près dans cette catégorie - \textit{sauf} dans le cas où $K_\infty$ est elle-même homotope à une catégorie
\msp{page 46}
finie. (Je ne me rends pas bien compte si ça arrive très souvent \ldots\ Mais on ne peut fonder une théorie générale des dérivateurs et des catégories de modèles `dérivables' sur une hypothèse de ce genre.)

Mais supposons que le foncteur $H$ envisagé se factorise même par Hot, i.e. qu'il transforme quasi-isomorphismes en bijections. On a vu que cela équivaut à la condition qu'il puisse se représenter par un système inductif sur $\mathbf{N}$ formé de catégories de Thomason-Kan. Il doit être facile de plus, si on y tient, de s'arranger pour que les foncteurs de transition soient des monomorphismes et pleinement fidèles, pour avoir une $\varinjlim$ sympathique. Ce sera bien sûr une catégorie de Thomason-Kan. Cette catégorie est-elle déterminée par le foncteur $H : \Cathot^{o} \longrightarrow \Ens$ à isomorphisme unique près dans Cat\,hot? Bien sûr, il n'est pas question qu'elle représente le foncteur $H$, sauf \texttt{[si]} celui-ci est bel et bien représentable, i.e. si le ind-objet $\uuK$ de Cat\,hot est essentiellement constant (ce qui ne doit pratiquement jamais être le cas). Notons déjà
\printedpage{30}
% head: A. Grothendieck : Dérivateurs, Ch. V, Édition des Universités de Montpellier II et Paris VII
% foot: Version 9 septembre 2001
\msp{page 47}
que si $u$ est un homomorphimse de ind-objets \textit{dans} Cat, ou, ce qui revient maintenant au même, un homomorphisme dans Cat % [sic] homomorphimse -> homomorphisme, transposed letters
\[
K_\infty \xrightarrow{u_\infty} K'_\infty
\]
tel que $\forall n$, $u_\infty(K_n) \subseteq K'_{n'}$ pour $n' = \tau(n)$ convenable, et si on suppose que l'homomorphisme de ind-objets induit un isomorphisme pour les foncteurs ind-représentables \textit{sur} Cat\,hot (par exemple, si $\uuK$ et $\uuK'$ ind-représentent le même foncteur $H$, et l'homomorphisme de ind-objets est celui qui correspond à $\mathrm{id}_H$), alors $u_\infty$ est forcément un isomorphisme \textit{dans} Hot, donc aussi (sauf erreur) dans Cat\,hot puisque $K_\infty$, $K'_\infty$ sont des catégories de Thomason-Kan. Mais ce qui n'est pas évident, c'est l'existence d'un homomorphisme $u_\infty$ comme \texttt{[ci-]}dessus, et s'il existe, qu'il soit unique (ou du moins canonique?) comme flèche de Cat\,hot.

Pour l'existence d'un $u_\infty$, on est ramené à prouver que si on a un diagramme à flèches pleines
\[
\begin{tikzcd}
K_n \arrow[r, "i_n", "\text{cofib.}"'] \arrow[d, "u_n"] & K_{n+1} \arrow[dr, dashed] & \\
K'_{\tau(n)} \arrow[rr, "\text{transition}"'] & & K'_N ,
\end{tikzcd}
\]
et si une flèche $u_{n+1}$ est donnée `en pointillés' \textit{dans} Cat\,hot, de façon à rendre le carré commutatif,
\msp{page 48}
alors on peut la réaliser par une flèche \textit{dans} Cat qui rende le diagramme commutatif pour de bon (dans Cat). On est ramené au même problème dans une situation
\[
\begin{tikzcd}
X \arrow[r, hook, shift right, "i"'] & Y , \arrow[l, dashed, shift right, "p"']
\end{tikzcd}
\]
où $X$ est une catégorie de Thomason, $i$ une cofibration de Thomason, $p$ une rétraction de $Y$ dans $X$ \textit{dans} Cat\,hot, on veut la réaliser par une vraie rétraction (\footnote[7]{\texttt{[un ?\ \ dans la marge]}}).

En admettant que ça soit le cas, donc que tout homomorphisme des ind-objets $\uuK$, $\uuK'$ \textit{dans} Cat\,hot se représente par un homomorphisme entre les ind-objets \textit{dans} Cat, donc aussi par un foncteur $K_\infty \longrightarrow K'_\infty$, reste la question si l'homomorphisme correspondant dans Cat\,hot est unique. Donc si on a deux foncteurs $f$, $g$ qui induisent le même homomorphisme sur les foncteurs ind-représentés sur Cat\,hot, si $f$ et $g$ sont homotopes. Donc

\textbf{Question}. Deux foncteurs $f, g : K_\infty \longrightarrow K'_\infty$ tels que pour tout $n \in \mathbf{N}$, $f_n$, $g_n$ qui s'en déduisent par restriction à $K_n$ sont homotopes (et même homotopes en tant que foncteurs de $K_n$ dans un $K'_N$ convenable), sont-ils homotopes? (\footnote[8]{$K$, $K'$ catégories de Thomason, les $K_n$, $K'_n$ aussi.})

\printedpage{31}
% head: A. Grothendieck : Dérivateurs, Ch. V, Édition des Universités de Montpellier II et Paris VII
% foot: Version 9 septembre 2001
\msp{page 49}
Je présume que non - qu'il peut arriver que la longueur minimum requise $l(n)$ pour un chemin dans $\uHom(K_n, K'_{\tau(n)})$ reliant $f_n$ et $g_n$, tende vers l'infini si $n \longrightarrow \infty$. L'hypothèse que les $K'_n$ sont des catégories de Thomason y change-t-il quelque chose (deux objets homotopes dans $\uHom(X, K')$, si $K'$ est de Thomason, sont-ils reliés par une flèche?)

Bien sûr, on peut prétendre que les catégories $\uCh_{\Sigma,n}(\mathcal{M}; a, b)$ ont été définies `canoniquement', et leur limite inductive $\uCh_{\Sigma}(\mathcal{M}; a, b)$ aussi. Mais la définition était assez arbitraire et abracadabrante - je suis convaincu que dix mathématiciens, se proposant de développer cette même idée des catégories de chemins, dans le même but que moi, arriveraient à dix définitions différentes! Pourtant, il n'est pas exclu que l'on finisse par trouver un système transitif d'isomorphismes dans Cat\,hot, entre les catégories de chemins `infinis' (i.e. de longueur non précisée) $K_\infty = K_\infty(\mathcal{M}; a, b)$ que nous aurons construites, dans le but
\msp{page 50}
principal (disons) de représenter le foncteur $H$ \textit{sur les catégories finies.} Et que les isomorphismes obtenus seront, bien entendu, des homomorphismes dans Ind(Cat\,hot). J'ai l'impression de ne pas avoir encore axiomatisé à fond la notion même de chemins, qui après tout a un sens géométrique assez évident et important par lui-même. Donc par ce qu'il faut entendre par `\textit{une} théorie des chemins dans une catégorie'. Une telle théorie $\mathcal{T}$ associerait donc à $\mathcal{M}$, $\Sigma$, $a$, $b$ un système inductif de catégories $\Ch^{\mathcal{T}}_{n,\Sigma}(\mathcal{M}; a, b)$. Et les contraintes géométriques naturelles à imposer à de telles théories, en vue aussi de leur utilisation pour expliciter
\[
\Hom_{\mathcal{M}(X)\Sigma_X^{-1}}(a_X, b_X) \xrightarrow{\ \sim\ } \pi_0\uHom(X, \underline{\underline{\underline{\mathrm{Ch}}}}^{\mathcal{T}}(\mathcal{M}; a, b)) % ?? uncertain: number of underlines under Ch -- triple or double
\]
donnerait un système transitif d'isomorphismes non seulement entre les ind-objets de Cat\,hot, mais peut-être aussi entre les $\Ch^{\mathcal{T}}_{\infty,\Sigma}(\mathcal{M}; a, b)$, pour des théories differentes. Si on dispose d'un tel énoncé, la catégorie $\Ch_{\infty,\Sigma}(\mathcal{M}; a, b)$, ou plus exactement ses images dans Cat\,hot et Hot, % [sic] differentes -> différentes, missing accent
\msp{page 51}
prendraient un sens convaincant.

Il y a d'autre part un autre point de vue, qui peut-être pourrait `canonifier' les catégories $\uCh_\infty$ en tant qu'objet de Cat\,hot et surtout de Hot (on passerait à Cat\,hot en appliquant le foncteur $\psi$, enveloppe de Thomason-Kan \ldots). En effet, on est intéressé à construire, en termes de $\mathcal{M}$, $\Sigma$, un `dérivateur' (ou un prédérivateur) $\mathbf{D}_{\mathcal{M},\Sigma}$ surtout sur le `domaine' des catégories finies. L'étendre à un domaine plus large est peut-être important dans un `deuxième souffle', mais peut être considéré comme du luxe dans le premier mouvement de fondation d'une théorie des dérivateurs. Ceci dit, on peut se poser des questions du type suivant.

Soit
\[
\Hot_0 \subset \Hot
\]
\printedpage{32}
% head: A. Grothendieck : Dérivateurs, Ch. V, Édition des Universités de Montpellier II et Paris VII
% foot: Version 9 septembre 2001
la sous-catégorie pleine de Hot, définie par les catégories finies (\footnote[9]{\textit{NB} Hot$_0$ est équivalente à une petite catégorie.}) - appelons les types d'homotopie `de type fini' (ou de \textit{présentation} finie?). On a un foncteur évident
\msp{page 52}
\[
\begin{aligned}
\Hot \quad &\longrightarrow \quad \Hot_0^{\wedge} = \uHom(\Hot_0^{o}, (\Ens)) \, ,\\
X \quad &\longmapsto \quad \tilde{X} = (\xi \longmapsto \Hom_{\Hot}(\xi, X)).
\end{aligned}
\]
Ce foncteur est \textit{conservatif} \texttt{[cela semble faux]}, puisque un homomorphisme $X \longrightarrow X'$ dans Hot est un isomorphisme s'il induit des isomorphismes pour les foncteurs $\tilde{X} \longrightarrow \tilde{X}'$, et même seulement pour leurs valeurs pour les sphères $S^n$ ($n \geq 0$). (Sauf erreur. Y a-t-il des canulars, parce que le critère d'équivalence faible par les $\pi_i$ ne s'applique pas directement, en l'absence de points bases. Ça ne doit pas être très sérieux \ldots\ \texttt{[cela semble en fait sérieux]}). Mais est-il peut-être même \textit{pleinement fidèle} (\footnote[10]{\texttt{[un ?\ \ dans la marge]}})? De sorte que Hot se réaliserait comme sous-catégorie pleine de $\Hot_0^{\wedge}$, contenant les foncteurs représentables?
\end{document}
